%% JAIR Manuscript - Clément Frassier
%% Apprentissage fédéré : principes et mise en œuvre
%% Détection du stress chez les infirmières hospitalières (Projet HEART)

\documentclass[manuscript, screen, review]{jair}

\setcopyright{cc}
\acmDOI{10.1613/jair.1.xxxxx}

\JAIRAE{Clément Frassier}
\JAIRTrack{Applications}
\acmVolume{4}
\acmArticle{6}
\acmMonth{8}
\acmYear{2026}

\RequirePackage[
  datamodel=acmdatamodel,
  style=numeric,
  backend=biber,
  giveninits=true,
  uniquename=init
  ]{biblatex}

\renewcommand*{\bibopenbracket}{(}
\renewcommand*{\bibclosebracket}{)}

\addbibresource{sample-base.bib}

\begin{document}

\title[Apprentissage Fédéré - Détection de Stress]{Apprentissage fédéré : principes et mise en œuvre}

\author{Clément Frassier}
\orcid{0009-0000-0000-0000}
\email{clement.frassier@gmail.com}
\affiliation{%
  \institution{UKM / UEL}
  \city{Kuala Lumpur}
  \country{Malaysia}
}

\begin{abstract}
{\bf Background:}
Les données physiologiques collectées par bracelet connecté en milieu hospitalier
constituent des données de santé à caractère personnel soumises au
RGPD~\cite{gdpr2016}. Les agréger sur un serveur central expose les patients et le
personnel à des risques majeurs de ré-identification, rendant les approches
centralisées classiques inadaptées à ce contexte.

{\bf Objectives:}
Nous visons à entraîner un modèle de détection automatique du stress sur des données
distribuées entre 15 infirmières, sans jamais transmettre ces données à un serveur
central, tout en fournissant une garantie formelle de confidentialité différentielle
et en produisant un modèle déployable sur bracelet embarqué (TinyML).

{\bf Methods:}
Nous proposons un pipeline d'apprentissage fédéré combinant FedProx ($\mu = 0{,}05$)
pour la hétérogénéité Non-IID, et DP-SGD via Opacus pour la confidentialité
différentielle $(\varepsilon, \delta)$-DP. Le modèle MLP Tiny (1~362 paramètres,
5,32~Ko) est entraîné sur 30 rounds de communication et quantifié dynamiquement en
INT8 pour le déploiement. Une recherche sur grille de 25 runs fédérés
(5 valeurs de $\sigma \in [0{,}6;\, 1{,}8]$, 5 graines indépendantes) mesure
empiriquement le compromis confidentialité/utilité, comparée à une baseline
centralisée non contrainte (même modèle, mêmes 5 graines, sans FL ni DP) qui sert de
plafond de référence théorique -- non déployable en pratique sur ce type de données
de santé au sens du RGPD, mais utile pour quantifier honnêtement le coût de la
confidentialité.

{\bf Results:}
Sur un test-set groupé (toutes infirmières confondues), le modèle fédéré atteint une
précision de $82{,}0 \pm 0{,}5\,\%$ en moyenne sur les 5 valeurs de $\sigma$, contre
$81{,}95 \pm 0{,}89\,\%$ pour le centralisé -- des précisions brutes statistiquement
indiscernables. Sur des métriques équilibrées par classe, le centralisé conserve un
avantage apparent ($57{,}8 \pm 1{,}2\,\%$ de balanced-accuracy contre $\approx 47{,}9\,\%$
pour le FL). Mais 11 des 15 infirmières ont un test-set mono-classe, où une prédiction
constante obtient mécaniquement $100\,\%$ -- ces clients dominent le résultat groupé et
en masquent la vraie signification. Restreinte aux 4 clients au test-set réellement
mixte, la mesure la plus honnête cliniquement, l'écart se réduit fortement
($45{,}7 \pm 1{,}5\,\%$ pour le centralisé contre $42{,}7 \pm 1{,}0\,\%$ pour le FL) et
le \textbf{centralisé lui-même s'effondre} : son AUC-ROC tombe à $0{,}267$ (sous le
niveau du hasard) et son seuil de décision optimal se replie sur la prédiction
constante, exactement la même signature d'absence de signal généralisable que celle
attribuée au FL seul dans une évaluation groupée naïve. Une ablation FedProx sans DP-SGD
confirme que la fédération (pas la confidentialité) explique l'essentiel de l'écart
résiduel, la confidentialité ajoutant un coût mesurable mais secondaire, plus visible sur
ces clients difficiles ($45{,}2\,\%$ sans DP contre $42{,}7\,\%$ avec DP) que sur la
métrique groupée. Deux architectures de personnalisation (FedPer, puis une combinaison
tête multi-couches + DP restreint à l'extracteur + SCAFFOLD) réduisent cet écart de
façon statistiquement significative sur les clients difficiles ($48{,}4\,\%$ puis
$59{,}9\,\%$ de balanced-accuracy) sans le faire disparaître.

{\bf Conclusions:}
Le résultat central de ce travail n'est pas que le FL coûte cher face au centralisé --
c'est que \textbf{l'hétérogénéité extrême entre infirmières limite la généralisation de
toute architecture testée}, centralisée comprise : sur les clients au profil
physiologique le plus atypique, même un modèle entraîné sans aucune contrainte de
confidentialité ni de décentralisation échoue à produire un signal exploitable. Le FL
n'ajoute qu'un désavantage marginal, mais réel et mesurable, par-dessus ce problème de
fond. Ceci reformule la question de recherche : il ne s'agit pas de départager FL et
centralisé sur la précision (le centralisé restant de toute façon juridiquement
inadapté à ce type de données de santé au regard du RGPD, ce qui fait du FL la seule
approche viable en pratique), mais de comprendre pourquoi l'hétérogénéité Non-IID casse
la généralisation et comment la personnalisation architecturale (FedPer, SCAFFOLD) peut
la réduire. Nos résultats quantifient honnêtement ce coût résiduel, montrent qu'il reste
stable quel que soit le niveau de confidentialité visé, et qu'il peut être
significativement atténué par une architecture personnalisée -- tout en conservant un
modèle déployable sur microcontrôleur (5,32~Ko).
\end{abstract}

\received{20 May 2026}
\received[accepted]{1 July 2026}

\maketitle

\section{Introduction}

L'apprentissage fédéré (\textit{Federated Learning}, FL) est un paradigme
d'apprentissage automatique distribué dans lequel un modèle global est entraîné de
manière collaborative par un ensemble de participants, appelés \textit{clients}, sans
que leurs données locales ne soient jamais transmises à un serveur
central~\cite{mcmahan2017communication}. Seuls les paramètres du modèle (poids et
biais) circulent sur le réseau, ce qui positionne le FL comme une réponse directe aux
exigences de minimisation des données du Règlement Général sur la Protection des
Données (RGPD, Art.~5)~\cite{gdpr2016} et de l'HIPAA~\cite{hipaa1996} en contexte
médical.

Pour ancrer cette présentation, nous utilisons comme fil rouge un cas d'usage concret
développé dans le cadre du projet HEART\footnote{Health, Engineering, AI,
Responsibility and TinyML : collaboration entre l'Universiti Kebangsaan Malaysia et
l'University of East London.} : la \textbf{détection automatique du stress de
15 infirmières hospitalières} à partir de signaux physiologiques collectés par bracelet
connecté. Ce scénario illustre pourquoi le FL est non seulement pertinent mais, en
pratique, la seule approche viable lorsque les données sont simultanément sensibles,
hétérogènes et distribuées sur des appareils contraints.

Ce travail part d'une question de comparaison (le FL peut-il approcher la performance
d'un modèle centralisé sur ce type de données ?) mais aboutit à un résultat plus
fondamental : \textbf{l'hétérogénéité physiologique extrême entre infirmières limite la
généralisation de toute architecture testée, centralisée y comprise}. Sur les
sous-groupes de patientes au profil le plus atypique, même un modèle entraîné sans
aucune contrainte de confidentialité ni de décentralisation échoue à produire un signal
exploitable (Section~\ref{subsec:why_cent_wins}). Ce constat déplace la question de
recherche : il ne s'agit pas de départager le FL et le centralisé sur la précision, mais
de comprendre pourquoi le Non-IID casse la généralisation dans ce contexte clinique et
dans quelle mesure des architectures de personnalisation (FedPer, SCAFFOLD) peuvent en
réduire l'effet -- une lecture qui fait de ce travail une étude de cas sur les limites
du FL en présence d'hétérogénéité extrême, plutôt qu'une simple comparaison
FL-vs-centralisé.

\section{Présentation du projet et méthodologie}
\label{sec:methodology}

Cette section décrit les composants de notre pipeline d'apprentissage fédéré, la nature
du jeu de données physiologiques, les choix d'architecture réseau et la modélisation
de la confidentialité.

\subsection{Le cycle d'apprentissage fédéré}
\label{subsec:fl_cycle}

Un cycle FL se déroule sur $T$ rounds de communication. À chaque round $t$, le serveur
diffuse les poids globaux courants $\mathbf{w}^{(t)}$ vers un sous-ensemble de clients.
Chaque client $k$ effectue $E$ epochs d'entraînement local sur ses propres données
$\mathcal{D}_k$, puis renvoie ses poids mis à jour $\mathbf{w}_k^{(t+1)}$. Le serveur
agrège ces contributions pour former le nouveau modèle global :

\begin{equation}
  \mathbf{w}^{(t+1)} = \sum_{k=1}^{K} \frac{|\mathcal{D}_k|}{\sum_j |\mathcal{D}_j|}
  \, \mathbf{w}_k^{(t+1)}
  \label{eq:fedavg}
\end{equation}

Cette règle d'agrégation, connue sous le nom de \textit{FedAvg}~\cite{mcmahan2017communication},
est la brique fondamentale sur laquelle s'appuient les stratégies plus avancées
utilisées dans ce travail.

\subsection{Jeu de données et prétraitement}
\label{subsec:dataset}

Nous utilisons le \textit{Dataset for Continuous Stress Monitoring of Hospital
Nurses}~\cite{nurse_stress_dataset}, collecté auprès de 15 infirmières équipées du
bracelet clinique Empatica~E4. Les signaux enregistrés sont l'activité électrodermale
(EDA), la fréquence cardiaque (HR), la température cutanée (TEMP) et l'accélération
triaxiale (X, Y, Z), chacun étant un proxy physiologique direct de la réponse du
système nerveux sympathique au stress.

Les signaux physiologiques constituent des données de santé à caractère personnel au
sens du RGPD~\cite{gdpr2016}. Les agréger sur un serveur central créerait un risque
majeur pour la vie privée (point de défaillance unique) et exposerait le système à
des attaques par ré-identification. De plus, ces signaux sont intrinsèquement
\textbf{hétérogènes} (Non-IID) : un rythme cardiaque de 95 bpm peut correspondre à
un état de repos pour une infirmière et à un état de stress aigu pour une autre.

Les séries temporelles brutes sont transformées en vecteurs de caractéristiques via
une fenêtre glissante de 60 échantillons avec un pas de 30 (50~\% de recouvrement) :

\begin{equation}
  \mathbf{x}_i = \bigl[\,\mu,\, \sigma,\, \min,\, \max\,\bigr]
  \;\text{calcul\'es sur chaque signal dans la fen\^etre}\; i
  \label{eq:features}
\end{equation}

Ce qui donne $6 \times 4 = 24$ caractéristiques par fenêtre. Le label de la fenêtre
est assigné par vote majoritaire sur les labels bruts qu'elle contient.

\paragraph{Split temporel avec gap anti-fuite.} Le recouvrement de 50~\% entre fenêtres
consécutives implique que la fenêtre $i$ et la fenêtre $i+1$ partagent 30 échantillons
bruts sur 60. Un découpage train/test par permutation aléatoire des fenêtres placerait
fréquemment $i$ en entraînement et sa voisine quasi-identique $i+1$ en test, permettant
au modèle de mémoriser une partie du contenu du test plutôt que d'apprendre un signal
généralisable. Ce travail retient à la place, pour chaque infirmière, un découpage
\textbf{chronologique} : les premières 80~\% des fenêtres
(dans l'ordre d'enregistrement) forment le train, les 20~\% suivantes forment le test,
séparées par un \textbf{gap de 2 fenêtres} (60 échantillons bruts, soit la taille
d'une fenêtre entière) qui élimine tout chevauchement possible entre la dernière
fenêtre de train et la première fenêtre de test. Cette même logique de découpage par
client, appliquée avant tout regroupement, est aussi ce qui garantit l'absence de
fuite lorsque les jeux de test de plusieurs infirmières sont ensuite concaténés
(Section~\ref{subsec:global_comparison}). La standardisation est réalisée
\textbf{localement} sur chaque client (moyenne nulle, variance unitaire calculée
uniquement sur les fenêtres d'entraînement du client, puis appliquée au train et au
test), ce qui préserve la frontière de confidentialité et évite toute fuite des
statistiques de test vers l'entraînement.

\paragraph{Recombinaison binaire du label.} Le dataset source encode en réalité
\textbf{trois} niveaux de stress : classe 0 = pas de stress (18,8~\% des lignes
brutes), classe 1 = stress faible (7,0~\%), classe 2 = stress élevé (74,2~\%), ce
niveau étant lui-même dérivé d'un score continu moyenné sur chaque session
d'enregistrement~\cite{nurse_stress_dataset}. Ce travail retient l'intégralité des
données en fusionnant les classes 1 et 2 en une seule classe positive ``stress''
(label $= \mathbb{1}[\text{label brut} > 0]$), ce qui correspond à la question
cliniquement pertinente pour une alerte temps réel sur bracelet (présence de stress,
quel que soit son niveau) et permet d'exploiter les 15 infirmières du jeu de données.
Au total, le jeu de données comprend \textbf{383 584 fenêtres} (306 886
entraînement / 76 698 test),
réparties de manière très hétérogène entre les 15 clients : certaines infirmières
n'ont quasiment aucun épisode de stress sur toute leur session, d'autres sont
quasiment stressées en continu. Cette hétérogénéité extrême, propre au protocole de
collecte, motive l'évaluation sur un test-set global regroupé (Section~\ref{sec:results})
plutôt que sur des moyennes par client, qui seraient dominées par cet artefact de
composition des classes plutôt que par la performance réelle du modèle.

\paragraph{Une conséquence méthodologique importante : la majorité des clients ont un
test-set mono-classe.} Cette hétérogénéité n'est pas qu'une observation qualitative :
\textbf{11 des 15 infirmières ont un test-set ne contenant qu'une seule des deux
classes} (split chronologique déterministe, donc identique quelle que soit la graine).
Sur ces 11 clients, une prédiction constante obtient mécaniquement $100\,\%$ de rappel,
\emph{indépendamment de la qualité du modèle}. Seuls \textbf{4 clients} (identifiés
ci-après comme les clients 4, 5, 8 et 12) ont un test-set réellement mixte, et sont donc
les seuls à porter un signal informatif sur la capacité de généralisation du modèle.
Conséquence directe pour la suite de ce document : \textbf{toute métrique regroupant les
15 clients est mécaniquement tirée vers le haut par ces 11 clients triviaux}, et ne
reflète qu'imparfaitement la performance clinique réelle. Chaque tableau de résultats de
ce document rapporte donc systématiquement deux chiffres côte à côte : la métrique
regroupée sur les 15 clients (comparable à la littérature FL standard, où cette
composition de classe est rarement examinée client par client) et une métrique
restreinte aux 4 clients informatifs (la mesure la plus honnête du pouvoir discriminant
réel). Les deux racontent une histoire différente -- souvent avec un écart de plusieurs
dizaines de points -- et c'est précisément cet écart qui constitue l'un des résultats de
ce travail (Section~\ref{subsec:why_cent_wins}).

\subsection{Architecture du modèle : MLP Tiny}
\label{subsec:model}

Le modèle retenu est un \textbf{MLP Tiny} (Multi-Layer Perceptron minimaliste) composé
de trois couches entièrement connectées :

\begin{equation}
  \text{Lin\'eaire}(24 \to 32) \xrightarrow{\text{ReLU}}
  \text{Lin\'eaire}(32 \to 16) \xrightarrow{\text{ReLU}}
  \text{Lin\'eaire}(16 \to 2)
  \label{eq:mlp}
\end{equation}

Ce modèle totalise environ \textbf{1 362 paramètres}, correspondant à une empreinte
mémoire de \textbf{5,32~Ko} en précision virgule flottante 32 bits (FP32). Cette
architecture est justifiée par deux facteurs méthodologiques :

\begin{itemize}
  \item \textbf{Nature tabulaire de l'entrée.} Le vecteur d'entrée est composé de
  24 caractéristiques statistiques déjà agrégées temporellement. Il n'y a ni structure
  spatiale ni dépendance temporelle longue distance à exploiter. Des modèles plus
  complexes (CNN ou RNN) n'apporteraient aucun gain de performance discriminant et
  introduiraient des paramètres inutiles.
  \item \textbf{Empreinte mémoire compatible avec les MCU contraints.} Avec 1~362
  paramètres, le modèle s'intègre facilement dans la SRAM très limitée (16--256~Ko)
  des microcontrôleurs bas de gamme (Cortex-M0/M4) ciblés pour le wearable compagnon.
\end{itemize}

\subsection{Quantification dynamique INT8}
\label{subsec:quantization}

En vue du déploiement sur microcontrôleur, nous appliquons une \textbf{quantification
dynamique post-entraînement} (PTQ) via l'API PyTorch~\cite{paszke2019pytorch} :

\begin{verbatim}
net_q = torch.quantization.quantize_dynamic(
    model_cpu, {nn.Linear}, dtype=torch.qint8
)
\end{verbatim}

Cette opération convertit les poids des couches \texttt{Linear} de FP32 (32 bits) vers
INT8 (8 bits), réduisant la taille du modèle de \textbf{5,32~Ko à environ 1,4~Ko}
(facteur de compression de $\times 3{,}8$). Les activations restent en virgule flottante
à l'exécution (quantification \textit{dynamique}), ce qui supprime le besoin d'un jeu
de calibration, contrairement à la quantification statique.

La quantification est évaluée empiriquement à chaque round d'évaluation fédéré. Cette
opération est strictement non-destructive : le modèle continue de s'entraîner en
précision FP32 tout au long des rounds. À chaque évaluation, une copie du modèle local
est temporairement quantifiée en INT8 pour mesurer l'\textbf{erreur de quantification}
(différence de précision FP32 vs INT8), puis écartée.

\subsection{Stratégie d'agrégation : FedProx}
\label{subsec:fedprox}

En distribution Non-IID, les fonctions de perte locales divergent au fil des époques
locales, provoquant une dérive des poids locaux (\textit{client drift}). La convergence
de FedAvg n'est alors plus garantie. Nous adressons ce défi avec FedProx plutôt que
d'autres alternatives pour deux raisons :

\begin{itemize}
  \item Premièrement, FedBN repose sur l'exclusion des couches BatchNorm de
  l'agrégation, or le MLP Tiny n'en contient aucune ; FedBN se réduirait donc à FedAvg
  dans notre configuration.
  \item Deuxièmement, FedProx ne nécessite aucun état supplémentaire à communiquer
  entre serveur et clients (contrairement à SCAFFOLD qui transmet des variables de
  contrôle), maintenant le coût de communication par round à exactement 5,32~Ko.
\end{itemize}

FedProx~\cite{li2020fedprox} corrige cette dérive en ajoutant un terme proximal à la
fonction objectif locale de chaque client. Chaque client résout :

\begin{equation}
  \min_{\mathbf{w}} \left\{ F_k(\mathbf{w}) + \frac{\mu}{2}
  \| \mathbf{w} - \mathbf{w}^{(t)}_{\text{global}} \|^2 \right\}
  \label{eq:fedprox_form}
\end{equation}

La valeur $\mu = 0{,}05$ a été retenue comme équilibre optimal : un $\mu$ trop faible
($\approx 0$) ramène aux instabilités de FedAvg, tandis qu'un $\mu$ trop élevé
($\geq 0{,}5$) étouffe l'adaptation locale. Le terme proximal est calculé sur
l'ensemble des paramètres :

\begin{equation}
  \text{terme\_proximal} = \sum_{\ell} \left\|
  \mathbf{w}_\ell - \mathbf{w}_{\ell,\,\text{global}} \right\|^2_2
  \label{eq:proxterm}
\end{equation}

\subsection{Confidentialité différentielle : DP-SGD via Opacus}
\label{subsec:dp}

Bien que le FL évite le partage des données brutes, les attaques par inversion de
gradient~\cite{geiping2020inverting} peuvent reconstruire les données à partir des
paramètres. La confidentialité différentielle (DP) apporte une garantie mathématique
formelle~\cite{dwork2014foundations}. Un mécanisme $\mathcal{M}$ est $(\varepsilon,
\delta)$-DP si, pour tout ensemble de résultats $S$ et toute paire de jeux de données
adjacents $\mathcal{D}$ et $\mathcal{D}'$ :

\begin{equation}
  \Pr[\mathcal{M}(\mathcal{D}) \in S] \;\leq\; e^\varepsilon \cdot
  \Pr[\mathcal{M}(\mathcal{D}') \in S] + \delta
  \label{eq:dp_form}
\end{equation}

L'algorithme DP-SGD~\cite{abadi2016deep}, implémenté via Opacus~\cite{opacus},
applique à chaque étape de gradient :

\begin{enumerate}
  \item \textbf{Clipping des gradients par exemple} à une norme $\ell_2$ maximale $C$
  (\texttt{max\_grad\_norm}) :
  $\tilde{g}_i = g_i \cdot \min\!\left(1,\, C / \|g_i\|_2\right)$.
  \item \textbf{Injection de bruit gaussien} :
  $\hat{g} = \frac{1}{B}\!\left(\sum_{i=1}^B \tilde{g}_i +
  \mathcal{N}\!\left(0,\, \sigma^2 C^2 \mathbf{I}\right)\right)$,
  où $\sigma$ est le multiplicateur de bruit.
\end{enumerate}

Le budget cumulé $\varepsilon$ est suivi au cours des rounds via l'accountant PRV
(\textit{Privacy Random Variable}) d'Opacus~\cite{opacus}, successeur numérique du
\textit{moments accountant} original~\cite{abadi2016deep} (voir
Section~\ref{sec:discussion} pour une discussion de son coût mémoire).

\subsection{Protocole expérimental et baselines}
\label{subsec:protocol}

Les simulations sont orchestrées via le framework Flower~\cite{beutel2020flower}, qui
simule les 15 clients FL séquentiellement sur un seul nœud de calcul. Nous conduisons
une recherche sur grille complète faisant varier à la fois le multiplicateur de bruit
et la graine aléatoire : $\sigma \in \{0{,}6;\, 0{,}8;\, 1{,}0;\, 1{,}4;\, 1{,}8\}$,
chacun répété sur \textbf{5 graines indépendantes} (42, 123, 7, 2024, 31415), pour un
total de \textbf{25 runs fédérés}. Les paramètres fixes sont reportés au
Tableau~\ref{tab:dp_params_justif}.

\begin{table}[h!]
  \centering
  \caption{Hyperparamètres de confidentialité différentielle maintenus fixes.}
  \label{tab:dp_params_justif}
  \begin{tabular}{llp{6cm}}
    \toprule
    Paramètre & Valeur & Justification \\
    \midrule
    \texttt{max\_grad\_norm} ($C$) & 1,0 & Valeur recommandée pour la
    classification ; non incluse dans la grille, le focus étant sur $\sigma$. \\
    \texttt{dp\_delta} ($\delta$) & $10^{-5}$ & Valeur fixée avant l'élargissement du
    jeu de données (Section~\ref{subsec:dataset}). Avec $n = 306\,886$ fenêtres
    d'entraînement, $1/n \approx 3{,}26 \times 10^{-6}$ : $\delta = 10^{-5}$ dépasse
    désormais légèrement ce seuil recommandé $\delta < 1/n$~\cite{dwork2014foundations}
    (d'un facteur $\approx 3$). Ceci est documenté comme limite en
    Section~\ref{subsec:synthesis} ; les valeurs de $\varepsilon$ rapportées restent
    correctement calculées pour ce $\delta$, mais un $\delta = 10^{-6}$ serait
    recommandé pour une future itération. \\
    \bottomrule
  \end{tabular}
\end{table}

Plutôt que de s'appuyer sur un run unique par configuration, chaque paire (graine,
$\sigma$) de la grille, ainsi que la baseline centralisée, a été répétée sur
\textbf{5 graines aléatoires indépendantes}. Cette pratique est motivée par deux
considérations : d'une part, DP-SGD injecte du bruit gaussien à chaque étape de
gradient, ce qui introduit une variabilité non négligeable entre runs même avec des
données et hyperparamètres identiques ; d'autre part, la moyenne et l'écart-type sur
plusieurs runs permettent de vérifier que les conclusions ne dépendent pas d'une
initialisation aléatoire particulière.

Pour mesurer le coût ou gain du FL, nous entraînons en parallèle une
\textbf{baseline centralisée} : le même MLP Tiny entraîné de façon classique sur
toutes les données regroupées (sans DP, sans cloisonnement fédéré, avec
standardisation globale), moyennée sur les 5 mêmes graines (5 runs au total).

Le Tableau~\ref{tab:exp_config} récapitule la configuration des deux conditions
expérimentales.

\begin{table}[h!]
  \centering
  \caption{Configuration des conditions expérimentales de comparaison.}
  \label{tab:exp_config}
  \begin{tabular}{lll}
    \toprule
    Paramètre & FL (FedProx+DP) & Centralisé (baseline) \\
    \midrule
    Rounds / Epochs  & 30 rounds, 2 ep./round & 30 epochs \\
    Clients          & 15 (1 par infirmière)  & aucun (données regroupées) \\
    Standardisation  & Locale (par client)    & Globale \\
    Confidentialité  & DP-SGD ($\sigma \in [0{,}6;\, 1{,}8]$, $C=1{,}0$) & Aucune \\
    Agrégation       & FedProx ($\mu=0{,}05$) & SGD classique \\
    Graines indép.   & 5                      & 5 \\
    Runs totaux      & 25 (5 graines $\times$ 5 $\sigma$) & 5 \\
    \bottomrule
  \end{tabular}
\end{table}

\section{Résultats expérimentaux}
\label{sec:results}

Cette section présente l'ensemble des mesures physiques, mémoires et statistiques
relevées lors de nos simulations. Aucun commentaire interprétatif n'est proposé ici.

\subsection{Comparaison globale : FL vs Centralisé}
\label{subsec:global_comparison}

Le jeu de données présente une hétérogénéité extrême de la classe ``stress'' entre
infirmières (Section~\ref{subsec:dataset}) : évaluer et moyenner une précision par
client, comme le ferait une évaluation FL classique, biaiserait fortement le résultat
vers les artefacts de composition de classe de chaque client plutôt que vers la
performance réelle du modèle. Nous évaluons donc chaque modèle final (le modèle
fédéré global non personnalisé pour le FL, le modèle unique pour le centralisé) sur
un \textbf{test-set unique regroupant les 15 infirmières} (76~698 fenêtres, 86,0~\%
de stress), construit en concaténant le test chronologique de chaque client
\emph{après} découpage individuel avec le gap anti-fuite -- exactement la même
méthodologie que celle déjà utilisée pour bâtir le test-set du centralisé, garantissant
une comparaison à armes égales. La précision brute étant peu informative sur un
jeu de test aussi déséquilibré (une prédiction constante ``stress'' obtiendrait déjà
86,0~\%), nous rapportons également la balanced-accuracy et le F1 macro. Mais la
balanced-accuracy groupée sur 15 clients hérite elle-même d'une limite : 11 de ces
clients ont un test-set mono-classe (Section~\ref{subsec:dataset}), où ils obtiennent
mécaniquement $100\,\%$ quel que soit le modèle. Nous rapportons donc systématiquement,
à côté de la métrique pooled, la même métrique restreinte aux 4 clients au test-set
réellement mixte (4, 5, 8, 12) -- la mesure la plus représentative de la performance
clinique réelle, discutée en détail Section~\ref{subsec:why_cent_wins}.
Le Tableau~\ref{tab:main_results} présente ces résultats, moyennés sur les 5 valeurs de
$\sigma$ pour le FL (chacune moyennée sur 5 graines) et sur 5 graines pour le
centralisé.

\begin{table}[h!]
  \centering
  \caption{Résultats globaux sur le test-set groupé (15 infirmières, 76~698 fenêtres) :
  FL (moyenne sur 5 $\sigma \times$ 5 graines) vs baseline centralisée (5 graines).}
  \label{tab:main_results}
  \begin{tabular}{lcc}
    \toprule
    Métrique & FL (FedProx + DP, tous $\sigma$) & Centralisé \\
    \midrule
    Précision (test-set groupé)  & $82{,}0 \pm 0{,}5\,\%$            & $81{,}95 \pm 0{,}89\,\%$ \\
    Balanced-accuracy (pooled, 15 clients)  & $47{,}9 \pm 0{,}2\,\%$ & $57{,}84 \pm 1{,}16\,\%$ \\
    \textbf{Balanced-accuracy (4 clients informatifs)}\textsuperscript{\ddag} & $42{,}73 \pm 0{,}97\,\%$ & $45{,}73 \pm 1{,}54\,\%$ \\
    F1 macro (pooled, 15 clients) & $45{,}4 \pm 0{,}2\,\%$           & $58{,}56 \pm 1{,}26\,\%$ \\
    \textbf{AUC-ROC (4 clients informatifs)}\textsuperscript{\ddag} & $0{,}2047 \pm 0{,}0232$ & $0{,}2669 \pm 0{,}0197$ \\
    Précision INT8 (interne, pondérée par client) & $81{,}62 \pm 0{,}17\,\%$ & $81{,}43 \pm 1{,}13\,\%$ \\
    Erreur de quantification (interne)  & $\approx -0{,}05\,\%$      & $0{,}52 \pm 0{,}65\,\%$ \\
    Taille modèle FP32          & 5,32~Ko                            & 5,32~Ko \\
    Communication par round     & 5,32~Ko                            & n/a \\
    Budget $\varepsilon$ ($\sigma=1{,}4$ recommandé) & $0{,}82 \pm 0{,}06$ & n/a \\
    RAM pic FL, entraînement seul / pendant \texttt{get\_epsilon()} (round 30) &
    $0{,}18 \pm 0{,}01$ / jusqu'à $33{,}65 \pm 3{,}30$~Mo\textsuperscript{†} & $10{,}88 \pm 0{,}00$~Mo \\
    \bottomrule
  \end{tabular}

  \vspace{2pt}
  {\footnotesize \textsuperscript{†}Le pic pendant \texttt{get\_epsilon()} varie avec
  $\sigma$ ; détail complet par valeur de $\sigma$ au Tableau~\ref{tab:ram_results}.
  On rapporte les deux valeurs ici (et non la seule valeur la plus favorable) pour que
  la comparaison au centralisé reflète le pic RAM réel, pas seulement le régime
  d'entraînement hors calcul de confidentialité.\\
  \textsuperscript{\ddag}11 des 15 clients ont un test-set mono-classe et faussent la
  métrique pooled (Section~\ref{subsec:dataset}) ; restreinte aux 4 clients au
  test-set réellement mixte (4, 5, 8, 12), c'est la mesure la plus honnête
  cliniquement -- voir la discussion complète Section~\ref{subsec:why_cent_wins}.}
\end{table}

\paragraph{Décomposition fédération vs confidentialité.} La comparaison ci-dessus
oppose "FL + DP" à "centralisé sans FL ni DP" -- deux conditions qui diffèrent
simultanément sur deux axes (décentralisation et bruit de confidentialité), sans les
isoler. Pour savoir laquelle des deux explique le coût mesuré, une troisième condition
a été ajoutée : FedProx fédéré, sans DP-SGD (ni clipping, ni bruit), même learning
rate que la condition FL+DP, volontairement non réajusté pour isoler l'effet du DP
toutes choses égales par ailleurs (le clipping de gradient modifie la taille effective
du pas d'optimisation ; ne pas recalibrer le LR compare bien le même optimiseur avec et
sans DP, pas deux réglages différents). Résultat sur le même test-set groupé, 5
graines :

\begin{table}[h!]
  \centering
  \caption{Décomposition du coût : centralisé (ni FL ni DP), FL sans DP (FedProx seul),
  FL + DP (moyenne sur 5 graines chacun, sauf FL+DP moyenné sur 5$\sigma \times$5
  graines).}
  \label{tab:dp_decomposition}
  \begin{tabular}{lccc}
    \toprule
    Métrique & Centralisé & FL sans DP & FL + DP \\
    \midrule
    Accuracy         & $81{,}95 \pm 0{,}89\,\%$ & $83{,}10 \pm 0{,}92\,\%$ & $82{,}00 \pm 0{,}62\,\%$ \\
    Balanced-accuracy (pooled, 15 clients) & $57{,}84 \pm 1{,}16\,\%$ & $48{,}68 \pm 0{,}73\,\%$ & $47{,}83 \pm 0{,}35\,\%$ \\
    \textbf{Balanced-accuracy (4 clients informatifs)}\textsuperscript{\ddag} & $45{,}73 \pm 1{,}54\,\%$ & $45{,}18 \pm 2{,}38\,\%$ & $42{,}73 \pm 0{,}97\,\%$ \\
    F1 macro (pooled, 15 clients) & $58{,}56 \pm 1{,}26\,\%$ & $46{,}10 \pm 0{,}75\,\%$ & $45{,}33 \pm 0{,}28\,\%$ \\
    \bottomrule
  \end{tabular}

  \vspace{2pt}
  {\footnotesize \textsuperscript{\ddag}Voir Section~\ref{subsec:dataset} : 11/15
  clients mono-classe faussent la métrique pooled.}
\end{table}

Sur la métrique pooled, FL sans DP (48,68~\%) et FL + DP (47,83~\%) restent à moins
d'un point l'un de l'autre -- un écart plus petit que leurs écarts-types respectifs.
Sur les 4 clients informatifs, l'écart est plus marqué (45,18~\% contre 42,73~\%,
soit 2,45 points) : \textbf{le bruit DP-SGD a un coût réel, mais secondaire}, et ce
coût est proportionnellement plus visible sur les clients les plus difficiles que sur
la métrique groupée qui le dilue. Dans les deux lectures, l'essentiel de l'écart avec
le centralisé est déjà présent sans aucun DP-SGD -- la fédération elle-même
(l'agrégation FedAvg/FedProx à travers 15 clients à la composition de classe
extrêmement hétérogène, Section~\ref{subsec:dataset}) domine largement le coût, la
confidentialité différentielle n'ajoutant qu'un supplément mesurable mais modeste.

La précision brute ne permet pas de départager les trois conditions (82,0~\% contre
83,1~\% contre 81,95~\%, dans le bruit statistique). Sur la métrique pooled, l'écart
semble net : le centralisé affiche un pouvoir discriminant apparent
(balanced-accuracy $57{,}84\,\%$, nettement au-dessus du hasard à $50\,\%$) contre un FL
proche voire légèrement en-dessous du hasard ($\approx 47{,}9\,\%$). Mais sur les 4
clients informatifs, \textbf{le centralisé s'effondre lui aussi} ($45{,}73\,\%$, à peine
au-dessus du hasard) -- un résultat discuté en détail Section~\ref{subsec:why_cent_wins},
qui change la nature de la conclusion : ce n'est plus "le centralisé résout le
problème que le FL ne résout pas", c'est "l'hétérogénéité extrême limite les deux
approches, le FL n'ajoutant qu'un désavantage marginal par-dessus" (voir
Section~\ref{sec:discussion}).

Pour donner une base chiffrée à ces deux lectures plutôt qu'un jugement visuel sur
moyenne $\pm$ écart-type, un test de Welch a été conduit sur les 25 checkpoints FL
(5 $\sigma \times$ 5 graines) contre les 5 checkpoints centralisés. Les 25 checkpoints
FL sont traités comme des réplicats d'une même condition -- ce qui est défendable
puisque la balanced-accuracy est quasi constante quel que soit $\sigma$
(Tableau~\ref{tab:epsilon_grid_extended}), et non comme 25 mesures indépendantes d'un
effet variant avec $\sigma$. Sur l'accuracy brute, l'écart n'est pas significatif
($t=0{,}11$, $p=0{,}92$ ; confirmé par un test de Mann-Whitney, $p=0{,}56$, plus robuste
à l'hypothèse de normalité fragile sur seulement 5 valeurs côté centralisé). Sur la
balanced-accuracy, l'écart est hautement significatif ($t=-19{,}2$, $p < 10^{-4}$ ;
Mann-Whitney $U=0$, séparation complète des deux groupes, $p < 10^{-4}$) : ce n'est pas
un artefact de variance d'échantillonnage.

\subsection{Compromis confidentialité vs utilité}
\label{subsec:privacy_utility}

Le Tableau~\ref{tab:epsilon_grid_extended} détaille l'évolution du budget de
confidentialité $\varepsilon$ accumulé au round 30, de la précision INT8 interne
(pondérée par client, à titre de suivi de convergence), et de la balanced-accuracy
sur le test-set groupé, en fonction du multiplicateur de bruit $\sigma$.

\begin{table}[h!]
  \centering
  \caption{Grille complète confidentialité/utilité (5 valeurs de $\sigma$, moyenne
  $\pm$ écart-type sur 5 graines, $\delta = 10^{-5}$, round 30).}
  \label{tab:epsilon_grid_extended}
  \begin{tabular}{cccc}
    \toprule
    $\sigma$ & $\varepsilon$ (round 30) & Précision INT8 (interne) & Balanced-acc. (groupée) \\
    \midrule
    0,6 & $5{,}29 \pm 0{,}30$ & $81{,}70 \pm 0{,}15\,\%$ & $47{,}87 \pm 0{,}47\,\%$ \\
    0,8 & $2{,}18 \pm 0{,}12$ & $81{,}86 \pm 0{,}14\,\%$ & $47{,}65 \pm 0{,}16\,\%$ \\
    1,0 & $1{,}32 \pm 0{,}13$ & $81{,}43 \pm 0{,}46\,\%$ & $47{,}85 \pm 0{,}24\,\%$ \\
    1,4 & $0{,}82 \pm 0{,}06$ & $81{,}50 \pm 0{,}43\,\%$ & $47{,}67 \pm 0{,}32\,\%$ \\
    1,8 & $0{,}54 \pm 0{,}04$ & $81{,}60 \pm 0{,}17\,\%$ & $48{,}12 \pm 0{,}38\,\%$ \\
    \bottomrule
  \end{tabular}
\end{table}

L'epsilon décroît de façon monotone et régulière avec $\sigma$, exactement comme
attendu mathématiquement, ce qui confirme le bon fonctionnement de l'accountant. La
précision INT8 et la balanced-accuracy groupée restent, elles, quasiment constantes
sur toute la plage -- le compromis confidentialité/utilité est empiriquement plat,
qu'on le lise sur la métrique interne ou sur la métrique groupée : le coût mesuré en
Section~\ref{subsec:global_comparison} ne s'aggrave pas lorsqu'on exige une
confidentialité plus forte.

\subsection{Ablation : impact de la personnalisation locale}
\label{subsec:ablation_personalization}

Le Tableau~\ref{tab:ablation_b} compare, pour chaque valeur de $\sigma$, la précision
\emph{interne} (moyenne non-pondérée sur les 15 clients, mesurée sur leur propre
test local) du modèle global reçu du serveur (avant adaptation) avec celle du même
modèle après les époques d'entraînement local FedProx de ce round.

\begin{table}[h!]
  \centering
  \caption{Écart local vs. global (moyenne $\pm$ écart-type sur 5 graines), précision
  interne par client, round 30.}
  \label{tab:ablation_b}
  \begin{tabular}{ccc}
    \toprule
    $\sigma$ & Modèle Global (interne) & Écart local $-$ global \\
    \midrule
    0,6 & $78{,}08 \pm 1{,}04\,\%$ & $+0{,}61 \pm 0{,}85$ pts \\
    0,8 & $77{,}47 \pm 0{,}36\,\%$ & $+1{,}12 \pm 0{,}47$ pts \\
    1,0 & $77{,}75 \pm 0{,}42\,\%$ & $+0{,}47 \pm 0{,}75$ pts \\
    1,4 & $77{,}48 \pm 0{,}43\,\%$ & $+0{,}81 \pm 0{,}85$ pts \\
    1,8 & $78{,}13 \pm 0{,}75\,\%$ & $+0{,}28 \pm 0{,}88$ pts \\
    \bottomrule
  \end{tabular}
\end{table}

La personnalisation locale n'apporte ici qu'un gain marginal et bruité (moins de 1,2
point en moyenne, parfois statistiquement nul compte tenu de l'écart-type). Ce
résultat indique que la valeur ajoutée de la personnalisation FedProx est faible sur
ce jeu de données : l'essentiel de l'écart de performance se joue entre FL et
centralisé (Section~\ref{subsec:global_comparison}), pas entre modèle global et
modèle personnalisé au sein du FL.

\subsection{Scénario d'un nouvel arrivant : adaptation d'un modèle pré-entraîné}
\label{subsec:late_joiner}

L'écart de personnalisation ci-dessus est mesuré sur des clients qui ont \emph{déjà}
participé à l'entraînement FL. Un scénario de déploiement plus réaliste, et plus
exigeant, est celui d'une infirmière qui rejoint la fédération \emph{après coup} : elle
reçoit le modèle global déjà entraîné sur les autres, sans jamais avoir contribué elle-même.
Pour le simuler, le client 8 (composition de classe équilibrée en train comme en test --
36,0~\% de stress en test, non mono-classe, contrairement à la majorité des autres
clients, Section~\ref{subsec:dataset}) est totalement exclu de l'entraînement fédéré
(30 rounds, $\sigma=0{,}8$, 3 graines) : sa contribution est neutralisée dans
l'agrégation FedProx (poids nul, \texttt{num-examples}$=0$), le modèle global final
n'a donc jamais vu ses données. Ce modèle est ensuite évalué directement sur les
données du client 8 ("jour 1", zéro adaptation), puis affiné sur ses seules données
d'entraînement locales (8 époques, taux d'apprentissage identique au reste du
pipeline) avant réévaluation.

\begin{table}[h!]
  \centering
  \caption{Scénario du nouvel arrivant (client 8, moyenne $\pm$ écart-type sur 3
  graines) : avant et après adaptation locale sur un modèle jamais entraîné sur ses
  données.}
  \label{tab:late_joiner}
  \begin{tabular}{lccc}
    \toprule
    Condition & Accuracy & Balanced-accuracy & F1 macro \\
    \midrule
    Jour 1 (zéro adaptation) & $36{,}56 \pm 0{,}92\,\%$ & $50{,}41 \pm 0{,}72\,\%$ & $27{,}41 \pm 1{,}57\,\%$ \\
    Après adaptation locale (8 époques) & $51{,}75 \pm 1{,}38\,\%$ & $\mathbf{58{,}51 \pm 1{,}00\,\%}$ & $\mathbf{51{,}43 \pm 1{,}51\,\%}$ \\
    \bottomrule
  \end{tabular}
\end{table}

Sans adaptation, le modèle global est à peine au niveau du hasard sur ce nouveau client
(balanced-accuracy $50{,}41\,\%$) -- il n'a jamais vu sa distribution locale, et le
Non-IID extrême de ce jeu de données (Section~\ref{subsec:dataset}) limite sa
généralisation directe. Après seulement 8 époques de fine-tuning local, la
balanced-accuracy grimpe à $58{,}51\,\%$ (comparable au centralisé,
Tableau~\ref{tab:main_results}) et le F1 macro quasiment double. C'est un résultat plus
net que l'écart de personnalisation \og en régime \fg{} du Tableau~\ref{tab:ablation_b}
: l'apport réel de l'adaptation locale se révèle surtout pour un participant
véritablement nouveau, pas pour un client déjà intégré au cycle d'entraînement.

\textbf{Limite importante} : ce résultat repose sur un seul client exclu (le client 8)
et 3 graines, contre 5 graines partout ailleurs dans ce rapport -- un choix assumé pour
tenir le budget de temps disponible, pas une omission. C'est un \textbf{cas illustratif},
pas un résultat généralisable à l'ensemble des 15 infirmières : rien ne garantit qu'un
autre client exclu, avec une composition de classe différente, montrerait un gain
d'adaptation de même ampleur. Une étude systématique répétant cette procédure sur
plusieurs clients serait nécessaire pour une conclusion générale.

\subsection{Rappel par classe et cas extrême : le client le plus atypique}
\label{subsec:late_joiner_recall}

La balanced-accuracy du Tableau~\ref{tab:late_joiner} peut masquer des profils
d'erreur très différents selon la classe : un modèle qui prédit presque toujours la
même classe peut afficher une balanced-accuracy proche de $50\,\%$ sans que ce soit un
signe d'absence de signal -- exactement le piège identifié
Section~\ref{subsec:why_cent_wins}. Pour vérifier ce point sur le scénario du nouvel
arrivant, et pousser le test jusqu'au cas le plus défavorable, nous répétons la même
procédure (exclusion totale de l'entraînement, 30 rounds, $\sigma=0{,}8$, 3 graines,
8 époques d'adaptation) sur le client 5 -- précisément le client identifié
Section~\ref{subsec:why_cent_wins} comme concentrant 99,9~\% des erreurs confiantes du
modèle FL+DP global.

\begin{table}[h!]
  \centering
  \caption{Rappel par classe (seuil par défaut $0{,}5$), avant et après adaptation
  locale -- plage min--max sur 3 graines.}
  \label{tab:late_joiner_recall}
  \begin{tabular}{lcc}
    \toprule
    Client / Condition & Rappel non-stress & Rappel stress \\
    \midrule
    Client 8, jour 1 (zéro adapt.)      & $0{,}00$--$2{,}48\,\%$   & $99{,}93$--$100{,}00\,\%$ \\
    Client 8, après adapt. (8 ép.)      & $32{,}95$--$37{,}03\,\%$ & $81{,}17$--$86{,}28\,\%$ \\
    \addlinespace
    Client 5, jour 1 (zéro adapt.)      & $68{,}00$--$96{,}00\,\%$ & $0{,}00$--$1{,}02\,\%$ \\
    Client 5, après adapt. (8 ép.)      & $100{,}00\,\%^{\text{a}}$ & $3{,}13$--$19{,}66\,\%$ \\
    \bottomrule
  \end{tabular}
  \vspace{2pt}
  {\footnotesize $^{\text{a}}$Valeur exacte et identique sur les 3 graines (25/25
  fenêtres non-stress correctement classées à chaque fois) -- non une plage arrondie.
  À prendre avec prudence : le test du client 5 ne compte que 25 fenêtres non-stress
  au total (sur 1688), un dénominateur trop faible pour une estimation de rappel
  fiable sur cette seule classe.}
\end{table}

Le rappel par classe révèle deux échecs symétriques et opposés au jour 1 : le client 8
voit son modèle prédire presque exclusivement ``stress'' (rappel stress
$\approx 100\,\%$, rappel non-stress $\approx 0\,\%$), tandis que le client 5 montre le
biais inverse (rappel non-stress $68$--$96\,\%$, rappel stress $\leq 1{,}02\,\%$) --
cohérent avec l'AUC proche ou sous $0{,}5$ déjà mesurée pour ce client. Dans les deux
cas, la balanced-accuracy proche de $50\,\%$ du jour 1 cachait un biais systématique
fort, pas une absence de signal.

L'adaptation locale corrige partiellement ces deux biais en sens inverse, mais de façon
incomplète pour le client 5 : à seuil inchangé ($0{,}5$), le rappel stress ne remonte
qu'à $3{,}13$--$19{,}66\,\%$ (moyenne $12{,}15\,\%$) -- le modèle continue de manquer la
grande majorité des épisodes de stress réels sur ce client, ce qui serait cliniquement
inacceptable pour un système d'alerte. Le client 8, à l'inverse, atteint un rappel
stress de $81{,}17$--$86{,}28\,\%$ après adaptation, un profil beaucoup plus utilisable.

\paragraph{Diagnostic complémentaire (pas un protocole de déploiement) : calibration du
seuil post-adaptation.} Le même balayage de seuil que Section~\ref{subsec:why_cent_wins}
appliqué aux modèles adaptés révèle, pour les deux clients, un potentiel de rappel
supplémentaire non exploité au seuil par défaut -- mais dans des directions opposées :
le seuil optimal descend à $0{,}001$ pour le client 5 (rappel stress porté à
$26{,}46$--$66{,}63\,\%$, moyenne $42{,}53\,\%$) et monte à $0{,}950$ pour le client 8
(rappel stress ramené à $71{,}93$--$75{,}98\,\%$ au profit du rappel non-stress, porté à
$59{,}05$--$59{,}65\,\%$). \textbf{Ce balayage sélectionne le seuil qui maximise la
métrique sur le test-set utilisé ensuite pour la rapporter : c'est un outil de
diagnostic légitime pour révéler un problème de calibration, pas un protocole de
déploiement validé} -- le chiffre de rappel qui en résulte est une borne supérieure
optimiste, pas une performance attendue en production. On note en outre que le seuil
optimal du client 5 atteint la borne inférieure du balayage ($0{,}001$), le même
artefact de bord déjà rencontré Section~\ref{subsec:why_cent_wins} : la valeur exacte
reste donc provisoire et n'est pas rapportée comme un résultat final. La divergence de
direction entre les deux clients (vers $0$ pour l'un, vers $1$ pour l'autre) est
néanmoins cohérente avec leurs biais opposés au jour 1, et \emph{suggère} qu'un seuil
calibré par client -- sur un split de validation dédié, distinct du test, jamais
réalisé dans ce travail -- pourrait apporter un gain réel en déploiement. Avec
seulement $N=2$ clients testés, ceci reste une piste à vérifier, pas une conclusion
établie.

\textbf{Limites} : ces résultats reposent sur deux clients exclus (8 et 5) et 3 graines
chacun, contre 5 graines partout ailleurs dans ce rapport -- un choix assumé pour tenir
le budget de temps disponible. Ce sont des \textbf{cas illustratifs}, pas un résultat
généralisable aux 15 infirmières.

\subsection{Le même scénario sous FedPer : un nouvel arrivant sans historique de tête}
\label{subsec:fedper_late_joiner}

Le scénario du nouvel arrivant est rejoué sous FedPer (Section~\ref{subsec:fedper}) sur
les deux mêmes clients (8 et 5, exclusion totale de l'entraînement, 30 rounds,
$\sigma=0{,}8$, 3 graines). \textbf{La définition de ``jour 1'' diffère nécessairement
entre les deux architectures, et ce n'est pas un artefact à corriger mais une
conséquence directe de leur conception} : sous FedProx (Section~\ref{subsec:late_joiner}),
le nouvel arrivant hérite d'une tête \emph{mal calibrée mais informée} (celle du modèle
global final) ; sous FedPer, il n'existe littéralement \textbf{aucune tête à hériter} --
\texttt{fc3} n'est jamais transmis -- donc le ``jour 1'' FedPer est évalué avec une tête
\emph{vierge, initialisée aléatoirement}, posée sur l'extracteur partagé (jamais
entraîné sur les données de ce client). Ce ne sont pas deux mesures du même phénomène
toutes choses égales par ailleurs : ne pas les lire comme directement comparables terme
à terme.

\begin{table}[h!]
  \centering
  \caption{FedPer, scénario du nouvel arrivant : jour 1 (tête vierge aléatoire) vs après
  onboarding (entraînement de \texttt{fc3} seul, extracteur gelé, 8 époques) -- plage
  min--max sur 3 graines.}
  \label{tab:fedper_late_joiner}
  \begin{tabular}{lcc}
    \toprule
    Client / Condition & Balanced-accuracy & Rappel (non-stress~/~stress) \\
    \midrule
    Client 8, jour 1 (tête vierge)   & $50{,}06$--$69{,}66\,\%$ & $0{,}12$--$93{,}76\,\%$~/~$40{,}58$--$100{,}00\,\%$ \\
    Client 8, après onboarding       & $54{,}80$--$69{,}57\,\%$ & $62{,}37$--$80{,}86\,\%$~/~$43{,}14$--$58{,}28\,\%$ \\
    \addlinespace
    Client 5, jour 1 (tête vierge)   & $5{,}31$--$50{,}00\,\%$  & $4{,}00$--$100{,}00\,\%$~/~$0{,}00$--$6{,}61\,\%$ \\
    Client 5, après onboarding       & $73{,}33$--$80{,}41\,\%$ & $96{,}00$--$100{,}00\,\%$~/~$46{,}66$--$64{,}82\,\%$ \\
    \bottomrule
  \end{tabular}
\end{table}

\paragraph{Une source de variance supplémentaire, distincte de celle déjà documentée.}
Le jour 1 FedPer présente une dispersion considérable entre graines (client 8 :
balanced-accuracy 50,1~\% à 69,7~\% ; client 5 : 5,3~\% à 50,0~\%), nettement plus large
que le jour 1 FedProx (déterministe à checkpoint donné). Cette variance a une cause
identifiée et distincte de la dérive temporelle déjà documentée pour le client 5
(Section~\ref{subsec:why_cent_wins}) : elle provient uniquement de l'initialisation
aléatoire de la tête vierge, une source de bruit qui n'existe pas sous FedProx où la
tête héritée est déterministe étant donné le modèle final. Pour le client 5, deux des
trois graines d'initialisation de tête s'effondrent vers une prédiction quasi constante
``non-stress'' (balanced-accuracy exactement 50~\%, rappel stress 0~\%) tandis que la
troisième produit un résultat pire que le hasard sur les deux classes simultanément
(5,3~\%) -- à ne pas confondre avec un signal de mauvaise qualité de l'extracteur, ces
deux échecs ont des mécanismes différents (un classifieur non entraîné, pas un
extracteur défaillant).

\paragraph{Après onboarding : un gain net et surtout beaucoup plus stable.} Entraîner
uniquement \texttt{fc3} pendant 8 époques (extracteur gelé) resserre fortement la
dispersion et améliore le rappel stress du client 5 au seuil par défaut ($0{,}5$, sans
aucun recalibrage de seuil) : $46{,}66$--$64{,}82\,\%$ (moyenne $57{,}19\,\%$), contre
$3{,}13$--$19{,}66\,\%$ (moyenne $12{,}15\,\%$) pour l'adaptation complète post-hoc de
FedProx au même seuil par défaut (Tableau~\ref{tab:late_joiner_recall}) -- et même
comparable au diagnostic de calibration de seuil de FedProx (non déployable en l'état,
Section~\ref{subsec:why_cent_wins}), obtenu ici \emph{sans} qu'aucun seuil n'ait été
retouché. Le client 8 montre un resserrement de la dispersion plutôt qu'un gain net
(balanced-accuracy moyenne comparable, écart-type divisé par plus de 2) : l'onboarding
FedPer stabilise surtout le résultat, l'effet est moins spectaculaire que pour le
client 5 mais dans la même direction.

\textbf{Limites} : mêmes réserves que pour le scénario FedProx (N=2 clients, 3 graines,
cas illustratif) -- avec ici une réserve supplémentaire propre à FedPer : la variance de
jour 1 dépendant de l'initialisation aléatoire de la tête, une graine d'initialisation
différente pourrait produire un jour 1 sensiblement différent, indépendamment de la
graine de partition des données.

\subsection{Quantification et dynamique des poids}

La plage dynamique des poids (définie par $\max(|w|) - \min(|w|)$) a été relevée sur
le modèle final après 30 rounds d'apprentissage. Le Tableau~\ref{tab:dynamic_range}
présente cette mesure sur l'ensemble du réseau (dynamique globale) ainsi que sur la
première couche linéaire (\texttt{fc1.weight}).

\begin{table}[h!]
  \centering
  \caption{Plage dynamique des poids (moyenne $\pm$ écart-type sur 5 graines) selon
  le mode d'entraînement et le niveau de bruit DP.}
  \label{tab:dynamic_range}
  \begin{tabular}{lcc}
    \toprule
    Configuration & Dynamique Globale & Dynamique \texttt{fc1.weight} \\
    \midrule
    Centralisé (std.\ globale)  & $9{,}4467 \pm 1{,}0336$ & $7{,}1129 \pm 0{,}5652$ \\
    FL + DP ($\sigma = 0{,}6$)  & $2{,}9734 \pm 0{,}2251$ & $0{,}8979 \pm 0{,}0547$ \\
    FL + DP ($\sigma = 0{,}8$)  & $3{,}2638 \pm 0{,}4708$ & $0{,}9624 \pm 0{,}0361$ \\
    FL + DP ($\sigma = 1{,}8$)  & $3{,}0734 \pm 0{,}3360$ & $1{,}4376 \pm 0{,}1326$ \\
    \bottomrule
  \end{tabular}
\end{table}

\subsection{Empreinte mémoire (RAM pic)}

L'empreinte mémoire vive (RAM) maximale a été mesurée dynamiquement sur un client au
round 30, juste avant et juste après l'appel à \texttt{get\_epsilon()} d'Opacus.

\begin{table}[h!]
  \centering
  \caption{Consommation RAM pic (Mo, moyenne $\pm$ écart-type sur 5 graines) au
  round 30, avant et pendant le calcul de la confidentialité, en fonction de $\sigma$.}
  \label{tab:ram_results}
  \begin{tabular}{cccc}
    \toprule
    $\sigma$ & Niveau de bruit & Avant \texttt{get\_epsilon()} & Pendant \texttt{get\_epsilon()} \\
    \midrule
    0,6 & Faible    & $0{,}18 \pm 0{,}00$~Mo & $33{,}65 \pm 3{,}30$~Mo \\
    0,8 & Modéré-faible & $0{,}18 \pm 0{,}00$~Mo & $20{,}30 \pm 1{,}48$~Mo \\
    1,0 & Modéré    & $0{,}18 \pm 0{,}00$~Mo & $16{,}05 \pm 1{,}70$~Mo \\
    1,4 & Fort      & $0{,}18 \pm 0{,}00$~Mo & $13{,}44 \pm 1{,}00$~Mo \\
    1,8 & Très fort & $0{,}37 \pm 0{,}42$~Mo$^{*}$ & $11{,}34 \pm 0{,}43$~Mo \\
    \bottomrule
  \end{tabular}

  \vspace{2pt}
  {\footnotesize $^{*}$4 des 5 graines donnent $\approx 0{,}18$~Mo, cohérent avec les
  autres lignes ; une seule graine produit une valeur isolée à $1{,}11$~Mo, tirant la
  moyenne vers le haut. Vraisemblablement un artefact ponctuel de mesure
  (\texttt{tracemalloc} / passage du ramasse-miettes Python) plutôt qu'un effet réel
  de $\sigma$, sans quoi la même hausse serait observée sur les autres graines.}
\end{table}

\section{Interprétation et discussion}
\label{sec:discussion}

Cette section propose une analyse théorique et méthodologique des phénomènes observés.

\subsection{Ce que révèle vraiment l'écart FL/centralisé : une limite partagée, pas un
avantage du centralisé}
\label{subsec:why_cent_wins}

Sur le test-set groupé (15 clients), le centralisé affiche un avantage apparent en
balanced-accuracy (57,84~\% contre $\approx$ 47,9~\% pour le FL,
Section~\ref{subsec:global_comparison}). \textbf{Cette lecture ne survit pas à
l'inspection par client.} 11 des 15 clients ont un test-set mono-classe
(Section~\ref{subsec:dataset}) où toute prédiction constante obtient mécaniquement
$100\,\%$ : ils dominent la métrique pooled et masquent ce qui se passe réellement sur
les 4 clients (4, 5, 8, 12) dont le test-set est effectivement mixte -- les seuls sur
lesquels une balanced-accuracy est une mesure signifiante du pouvoir discriminant du
modèle.

\paragraph{Sur les clients informatifs, le centralisé s'effondre aussi.} Restreinte à
ces 4 clients, la balanced-accuracy du centralisé tombe à $45{,}73 \pm 1{,}54\,\%$ --
à peine au-dessus du hasard, et son AUC-ROC à $0{,}2669 \pm 0{,}0197$. \textbf{Prudence
sur cet effectif} : cette AUC groupée repose sur seulement 4 clients, un effectif
restreint du même ordre de grandeur que celui pour lequel ce document applique déjà une
réserve explicite (les 25 fenêtres non-stress du client 5 seul,
Tableau~\ref{tab:late_joiner_recall}) -- une partie de l'écart à $0{,}5$ pourrait relever
du bruit d'échantillonnage plutôt que d'une relation systématiquement apprise à l'envers,
même si la cohérence de signe avec le diagnostic par client ci-dessous rend cette lecture
peu probable. Par la même logique déjà établie pour le diagnostic du client 5 plus bas
dans cette section :
\textbf{une AUC significativement inférieure à $0{,}5$ ne signale pas une absence de
signal, mais un classement partiellement \emph{inversé}} -- les probabilités prédites
par le modèle centralisé classent les fenêtres de test dans le sens contraire de la
vérité terrain plus souvent que ne le ferait un classement aléatoire. Le centralisé n'a
donc pas simplement ``rien appris'' sur ces clients ; il y a appris une relation
partiellement fausse. Le balayage de seuil de décision (de $10^{-3}$ à $0{,}99$, sur les
probabilités déjà produites, aucun réentraînement) le confirme indirectement : le seuil
optimal du centralisé se replie sur $0{,}001$ (prédire ``stress'' quasi systématiquement,
soit inverser la prédiction par défaut plutôt que la calibrer), plafonnant à
\emph{exactement} $50{,}00 \pm 0{,}00\,\%$ -- une correction de seuil qui compense
l'inversion sans la résoudre, identique dans sa forme à ce que ce document attribuait
jusqu'ici au FL seul. Le FL+DP, sur ces mêmes 4 clients, présente la même signature,
un peu plus marquée (balanced-accuracy $42{,}73 \pm 0{,}97\,\%$, AUC
$0{,}2047 \pm 0{,}0232$, meilleur seuil également $0{,}001$, plafond à
$46{,}98 \pm 1{,}28\,\%$ -- sans même atteindre le plancher théorique, signe que
l'inversion y est moins complètement compensable par le seul choix du seuil).
\textbf{Le message central de ce document change en conséquence} : il ne
s'agit pas de ``pourquoi le centralisé résout un problème que le FL ne résout pas'',
mais de ``pourquoi aucune des deux architectures n'apprend une relation fiable sur les
clients les plus atypiques'' -- l'écart réel entre les deux (3 à 5 points de
balanced-accuracy selon la métrique) est un désavantage marginal du FL par-dessus un
problème de fond commun aux
deux, pas la preuve d'un FL structurellement inadapté.

\paragraph{Significativité statistique de l'écart centralisé vs FL, sur les clients
informatifs.} Un test de Welch et un test de Mann-Whitney ont été conduits sur la
balanced-accuracy des 4 clients informatifs : centralisé (5 graines) contre FL+DP
(25 valeurs, $5\sigma \times$ 5 graines), et centralisé contre FL sans DP (5 graines).
Résultat, cohérent avec la décomposition qui suit : \textbf{centralisé vs FL+DP est
statistiquement significatif} (Welch $t=4{,}21$, $p=9{,}8\times10^{-3}$ ; Mann-Whitney
$U=123/125$, $p=5{,}6\times10^{-5}$), mais \textbf{centralisé vs FL sans DP ne l'est
pas} (Welch $t=0{,}44$, $p=0{,}68$ ; Mann-Whitney $U=15/25$, $p=0{,}69$) -- sur ces
clients précis, la fédération seule (sans confidentialité) n'est pas statistiquement
distinguable du centralisé, et c'est l'ajout du bruit DP-SGD qui crée l'écart mesurable
avec le centralisé. Une nuance par rapport à la métrique pooled, où l'essentiel de
l'écart était attribué à la décentralisation plutôt qu'à la confidentialité -- cohérent
avec le constat déjà fait que le coût du bruit DP est proportionnellement plus visible
sur les clients difficiles que sur la métrique groupée.

\textbf{Vérification que ce contraste (significatif contre non significatif) ne
reflète pas simplement un déséquilibre de puissance statistique} entre les deux
comparaisons (FL+DP : $n=25$ contre FL sans DP : $n=5$, la même taille que le
centralisé) -- un piège classique déjà démonté ailleurs dans ce document (Section
\ref{subsec:duration_vs_granularity}). Deux vérifications indépendantes : (1)~la
taille d'effet (Cohen's $d$), insensible à ce déséquilibre, reste très supérieure pour
centralisé vs FL+DP ($d=2{,}81$, effet large) que pour centralisé vs FL sans DP
($d=0{,}28$, effet petit) -- un facteur 10 entre les deux tailles d'effet, pas
seulement entre les deux $p$-values ; (2)~un sous-échantillonnage répété de FL+DP à
$n=5$ (2\,000 tirages sans remise, pour égaler exactement la puissance du test
centralisé vs FL sans DP) reste significatif à $p<0{,}05$ dans $99{,}7\,\%$ des tirages
(p médian $=9{,}2\times10^{-3}$, quasi identique au $p=9{,}8\times10^{-3}$ sur
l'échantillon complet). Le contraste observé n'est donc pas un artefact de taille
d'échantillon : à puissance égale, l'écart centralisé/FL+DP reste presque toujours
détecté, contrairement à l'écart centralisé/FL sans DP.

\textbf{Limite à ne pas occulter sur la portée de cette significativité.} Les 25 (ou 5)
valeurs testées proviennent de 5 graines d'entraînement et 5 valeurs de $\sigma$ -- des
sources de bruit d'entraînement, pas de diversité clinique. La significativité mesurée
démontre que l'écart est \emph{reproductible d'une graine et d'un niveau de bruit à
l'autre, sur ces 4 clients précis} ; elle ne démontre pas que le phénomène se
généraliserait à un 5\textsuperscript{e} profil de client au test-set mixte, puisque ce
jeu de données n'en fournit aucun autre pour le vérifier. C'est une caractéristique
structurelle de ce dataset (seulement 4 clients au test-set réellement mixte,
Section~\ref{subsec:dataset}), pas une faiblesse de la méthode statistique employée --
mais elle borne strictement la portée de la conclusion : un résultat solide sur
\emph{ces} 4 infirmières, pas encore une généralisation démontrée à toute population de
patientes hétérogènes.

\paragraph{Décomposition fédération vs confidentialité, sur les deux métriques.}
Le Tableau~\ref{tab:dp_decomposition} isole la contribution du bruit DP-SGD :
\begin{itemize}
  \item \textbf{Sur la métrique pooled}, FL sans DP (48,68~\%) et FL + DP (47,83~\%)
  sont à moins d'un point l'un de l'autre -- le bruit DP-SGD y semble quasi gratuit.
  \item \textbf{Sur les 4 clients informatifs}, l'écart est plus net (FL sans DP
  $45{,}18 \pm 2{,}38\,\%$ contre FL+DP $42{,}73 \pm 0{,}97\,\%$, soit 2,45 points) :
  le coût du bruit de confidentialité est réel, simplement dilué par la métrique
  pooled qui le noie sous les clients triviaux. Il reste néanmoins secondaire face à
  l'écart avec le centralisé, qui persiste presque à l'identique avec ou sans DP-SGD.
\end{itemize}
Dans les deux lectures, l'agrégation FedAvg/FedProx à travers 15 clients aussi
hétérogènes (Section~\ref{subsec:dataset}) domine le coût mesuré ; la confidentialité
différentielle y ajoute un supplément mesurable mais net à moins que ne le suggère la
métrique pooled.

\paragraph{Concentration des erreurs, revisitée.} Une inspection par client de l'AUC
pooled avait révélé que 99,9~\% des fenêtres où le modèle FL+DP est confiant à tort
proviennent d'un seul client (le client 5) -- interprété comme un cas de
non-stationnarité Non-IID localisée plutôt qu'une inversion systémique de l'agrégation.
Ce constat se lit maintenant différemment à la lumière de ce qui précède : le client 5
est l'un des 4 seuls clients à porter un signal informatif du tout -- il n'est donc pas
surprenant que les erreurs du modèle s'y concentrent, puisque c'est justement l'un des
rares endroits où il y a quelque chose de non trivial à réussir ou à rater. Le
diagnostic (non-stationnarité Non-IID entre clients, pas un bug d'agrégation) reste
valable, mais s'inscrit maintenant dans un problème plus large touchant aussi le
centralisé, pas une pathologie propre au FL.

\paragraph{Une mitigation testée : loss pondérée par classe.} Puisque le déséquilibre
de classe par client (Section~\ref{subsec:dataset}) semble au cœur du problème, une
piste naturelle est de pondérer la fonction de perte locale par l'inverse de la
fréquence globale des classes (calculée une fois sur l'agrégat des partitions
d'entraînement de tous les clients). Vérification ponctuelle à $\sigma=0{,}8$
(5 graines, pas l'ensemble de la grille) :

\begin{center}
\begin{tabular}{lccc}
  \toprule
  Condition ($\sigma=0{,}8$) & Accuracy & Balanced-accuracy & F1 macro \\
  \midrule
  FL + DP standard        & $81{,}84 \pm 0{,}49\,\%$ & $47{,}77 \pm 0{,}29\,\%$ & $45{,}33 \pm 0{,}19\,\%$ \\
  FL + DP + loss pondérée & $81{,}73 \pm 0{,}65\,\%$ & $47{,}66 \pm 0{,}40\,\%$ & $45{,}21 \pm 0{,}25\,\%$ \\
  \bottomrule
\end{tabular}
\end{center}

Résultat honnête et négatif : l'écart (0,11 point de balanced-accuracy) est dans le
bruit, la loss pondérée n'apporte aucun gain mesurable ici. C'est cohérent avec la
décomposition ci-dessus -- repondérer la perte \emph{locale} de chaque client ne peut
pas corriger un effet qui émerge de l'\emph{agrégation} entre 15 mises à jour aussi
hétérogènes (dont plusieurs clients mono-classe, Section~\ref{subsec:dataset}) : le
problème ne se situe pas au niveau de la fonction de perte individuelle. Des pistes
agissant directement sur l'agrégation (repondérer la contribution de chaque client à
la moyenne FedProx, ou une architecture de personnalisation plus poussée) semblent plus
prometteuses et restent à explorer.

Ce travail ne quantifie donc pas tant ``le prix de la confidentialité et de la
décentralisation'' que \textbf{la difficulté commune, pour toute architecture, à
généraliser sur une population d'infirmières aussi hétérogène} : même le centralisé,
qui ne subit aucune des contraintes du FL, échoue sur les clients les plus atypiques
(balanced-accuracy à peine au-dessus du hasard, AUC en dessous). Le désavantage
supplémentaire du FL par rapport au centralisé (3 à 5 points selon la métrique) reste
réel, mesurable et non négligeable -- mais c'est un supplément à un problème de fond,
pas le problème lui-même. Cette distinction compte pour l'interprétation : le centralisé
n'étant de toute façon pas une alternative réellement disponible pour ce cas d'usage
(agréger les signaux physiologiques bruts de 15 infirmières sur un serveur unique
constituerait un traitement de données de santé difficilement compatible avec le
principe de minimisation des données du RGPD, Art.~5~\cite{gdpr2016}, tel que rappelé
en introduction), le FL reste la seule approche viable en pratique -- et ce travail
montre qu'il n'ajoute qu'un coût marginal par-dessus une limite que le centralisé ne
résoudrait pas non plus. Les travaux futurs les plus prometteurs (features plus
informatives, architectures de personnalisation comme FedPer et FedRep+SCAFFOLD ci-après)
visent donc moins à ``rattraper le centralisé'' qu'à réduire l'effet de l'hétérogénéité
Non-IID elle-même, un problème qui dépasse le choix FL-vs-centralisé.

\subsection{Une architecture testée pour la personnalisation renforcée : FedPer}
\label{subsec:fedper}

La piste de ``personnalisation renforcée'' évoquée ci-dessus a été testée directement,
via \textbf{FedPer}~\cite{arivazhagan2019federated} : l'extracteur de caractéristiques
(\texttt{fc1}, \texttt{fc2}) reste agrégé globalement comme en FedProx, mais la couche
de classification finale (\texttt{fc3}) n'est \emph{jamais} transmise au serveur ni
agrégée -- elle reste locale à chaque client, entraînée en continu sur ses propres
données pendant les 30 rounds. Contrairement à Ditto (deux réseaux complets par client,
régularisés l'un vers l'autre par un terme proximal), FedPer ne modifie qu'un seul
réseau et ne nécessite qu'un filtre sur les clés du dictionnaire de poids transmis --
une différence d'implémentation qui s'est révélée pertinente : l'historique de ce dépôt
contient un bug de synchronisation documenté sur les branches Ditto (`f2d0b15`,
`a3d8107`, `6e6ee40` -- poids locaux réécrasés par erreur avec les poids globaux à
chaque round), une classe d'erreur qui n'existe structurellement pas avec FedPer
puisqu'il n'y a rien à resynchroniser.

\paragraph{Méthodologie d'évaluation -- une différence à assumer explicitement.}
Contrairement à FedProx (Section~\ref{subsec:global_comparison}), FedPer ne produit
\textbf{aucun modèle global unique évaluable seul} : le checkpoint serveur final ne
contient que l'extracteur (\texttt{fc1}, \texttt{fc2}), jamais de tête de classification.
Chaque fenêtre de test est donc scorée par le modèle propre au client dont elle provient
(extracteur partagé + tête locale de ce client), puis les prédictions de tous les
clients sont regroupées pour obtenir les métriques globales. Ce choix n'est pas un
contournement : c'est la seule évaluation cohérente avec l'architecture, et elle est en
fait \textbf{plus représentative d'un déploiement réel} que le modèle unique de FedProx
-- chaque wearable exécuterait de toute façon son propre modèle.

\begin{table}[h!]
  \centering
  \caption{FedPer vs FedProx+DP (méthodologie d'évaluation propre à chaque
  architecture, voir texte) -- moyenne sur 5$\sigma\times$5 graines (n=25 chacun).}
  \label{tab:fedper_comparison}
  \begin{tabular}{lcc}
    \toprule
    Métrique & FedProx + DP & FedPer \\
    \midrule
    Accuracy (pooled, 15 clients)      & $82{,}02 \pm 0{,}62\,\%$ & $83{,}21 \pm 1{,}65\,\%$ \\
    Balanced-accuracy (pooled, 15 clients) & $47{,}86 \pm 0{,}35\,\%$ & $65{,}36 \pm 2{,}14\,\%$ \\
    \textbf{Balanced-accuracy (4 clients informatifs)}\textsuperscript{\ddag} & $42{,}73 \pm 0{,}97\,\%$ & $\mathbf{48{,}44 \pm 3{,}65\,\%}$ \\
    \textbf{AUC-ROC (4 clients informatifs)}\textsuperscript{\ddag} & $0{,}2047 \pm 0{,}0232$ & $\mathbf{0{,}4417 \pm 0{,}1022}$ \\
    Communication (FIT, moyenne) & $5{,}320$~Ko          & $5{,}188$~Ko ($-2{,}5\,\%$) \\
    \bottomrule
  \end{tabular}

  \vspace{2pt}
  {\footnotesize \textsuperscript{\ddag}Métrique restreinte aux clients au test-set
  réellement mixte (Section~\ref{subsec:dataset}) -- la mesure la plus honnête
  cliniquement.}
\end{table}

Sur la métrique pooled, l'écart est très marqué (Welch $t=40{,}36$,
$p=1{,}7\times10^{-24}$ ; Mann-Whitney $U=625/625$, séparation complète) -- mais ce
chiffre est gonflé par les 11 clients mono-classe (Section~\ref{subsec:dataset}), sur
lesquels les deux architectures obtiennent des scores très proches et souvent
identiques. \textbf{Sur les 4 clients informatifs, l'écart réel est plus modeste mais
reste statistiquement solide} (Welch $t=7{,}56$, $p=3{,}6\times10^{-8}$ ;
Mann-Whitney $U=590/625$, $p=7{,}7\times10^{-8}$) : +5,7 points de balanced-accuracy en
faveur de FedPer, avec une variance nettement plus grande que sur la métrique pooled
(écart-type $3{,}65\,\%$ contre $2{,}14\,\%$) -- FedPer aide sur les clients difficiles,
mais moins spectaculairement, et moins uniformément, que ne le suggère la métrique
pooled. Le détail complet (résultats par $\sigma$ et par graine, tests statistiques sur
les deux métriques) est conservé dans \texttt{resultsfeat/fedper\_results\_summary.md},
séparément de ce document.

\paragraph{Interprétation -- une confirmation directe, mais d'ampleur plus mesurée
qu'il n'y paraît.} Ce résultat complète la décomposition FL-vs-DP ci-dessus : celle-ci
avait attribué le coût de balanced-accuracy à l'\emph{agrégation}, pas au bruit de
confidentialité. FedPer teste directement cette hypothèse en supprimant l'agrégation de
la seule couche responsable de la décision finale -- et une partie du coût disparaît,
sur la métrique qui compte cliniquement. L'écart mesuré sur les clients informatifs
(vérifié également en écartant l'hypothèse d'un mélange accidentel de checkpoints entre
les deux conditions) ne doit pas surprendre, mais \textbf{ne doit pas non plus être
imputé à un facteur unique} : FedPer et FedProx diffèrent ici simultanément sur \emph{deux axes}
-- (1)~la \textbf{durée} de personnalisation (30 rounds continus pendant l'entraînement,
contre 8 époques de fine-tuning post-hoc pour FedProx, Tableau~\ref{tab:ablation_b}) et
(2)~la \textbf{granularité} (seule la tête de classification pour FedPer, contre le
réseau entier pour le fine-tuning post-hoc de FedProx). \textbf{Ces deux facteurs ont
depuis été isolés séparément} (Section~\ref{subsec:duration_vs_granularity}, à partir
du même checkpoint de départ et avec un fine-tuning re-seedé pour un test statistique
apparié sur 5 graines) : la \emph{granularité} domine largement (+6,5 à +7,4 points de
balanced-accuracy, $p<0{,}03$) -- personnaliser uniquement la tête, à partir du même
extracteur, surpasse nettement le fine-tuning du réseau entier, probablement en
limitant le surapprentissage sur le peu de données locales disponibles par client -- et
la \emph{durée} compte aussi, mais beaucoup plus modestement (+1 à +2 points,
$p<0{,}02$).
Fait notable, l'accuracy brute progresse elle aussi légèrement ($+1{,}19$ point,
$p=1{,}97\times10^{-3}$) : FedPer n'échange donc pas de la précision brute contre du
rappel équilibré, les deux progressent ensemble ici.

\subsection{Durée ou granularité ? La granularité domine, la durée compte un peu}
\label{subsec:duration_vs_granularity}

Pour trancher lequel des deux facteurs ci-dessus domine, FedProx+DP ($\sigma=0{,}8$,
checkpoint final) est fine-tuné hors ligne sur les données d'entraînement locales de
chacun des 4 clients informatifs, en repartant du \textbf{même} checkpoint dans toutes
les conditions et en ne faisant varier qu'un facteur à la fois -- durée (8 contre 30
époques) et granularité (réseau entier dégelé contre tête \texttt{fc3} seule dégelée).
Un premier test comparait directement le fine-tuning complet à FedPer (Tier 1) : cette
comparaison confondait la granularité avec le fait que FedPer part de \emph{son propre}
extracteur, entraîné selon un régime d'agrégation entièrement différent (30 rounds sans
terme proximal, jamais influencé par un \texttt{fc3} agrégé) -- corrigé ci-dessous en
repartant systématiquement du même point de départ. Le fine-tuning est explicitement
re-seedé à chaque graine pour permettre un test statistique apparié rigoureux (5 valeurs
par condition, comme partout ailleurs dans ce document).

\begin{center}
\begin{tabular}{lcc}
  \toprule
  Condition (même checkpoint de départ) & 8 époques & 30 époques \\
  \midrule
  Réseau entier dégelé      & $48{,}49 \pm 1{,}33\,\%$ & $50{,}43 \pm 1{,}71\,\%$ \\
  \textbf{Tête seule (\texttt{fc3}) dégelée} & $\mathbf{55{,}88 \pm 2{,}35\,\%}$ & $\mathbf{56{,}93 \pm 2{,}66\,\%}$ \\
  \bottomrule
\end{tabular}
\end{center}
{\footnotesize Balanced-accuracy (moyenne $\pm$ écart-type, 5 graines), pooled sur les
4 clients informatifs. Référence : FedProx+DP non adapté $42{,}73 \pm 0{,}97\,\%$ ;
FedPer (son propre extracteur, 30 rounds continus) $48{,}44 \pm 3{,}65\,\%$.}

Un test $t$ apparié (5 graines, chaque graine comparée à elle-même entre conditions)
quantifie les deux effets : \textbf{la granularité domine largement}
($+7{,}39$ points à 8 époques, $t=5{,}26$, $p=6{,}2\times10^{-3}$ ; $+6{,}51$ points à
30 époques, $t=3{,}64$, $p=2{,}2\times10^{-2}$), \textbf{mais la durée compte aussi,
modestement} ($+1{,}94$ point pour le réseau entier, $t=8{,}11$,
$p=1{,}3\times10^{-3}$ ; $+1{,}05$ point pour la tête seule, $t=4{,}43$,
$p=1{,}2\times10^{-2}$) -- \textbf{une révision par rapport à un premier test non
seedé qui avait conclu, à tort, que la durée n'avait aucun effet} : l'absence de
contrôle du bruit d'ordre de mélange masquait un effet réel mais petit, cohérent en
signe sur les 5 graines une fois ce bruit maîtrisé. Le test de Wilcoxon apparié, en
complément, plafonne à $p=0{,}0625$ dans les quatre cas -- un artefact de plancher
propre à $n=5$ paires dont toutes les différences ont le même signe (le $p$ minimal
atteignable avec si peu de graines), pas un signe de significativité douteuse ; le
test $t$ reste la lecture appropriée ici.

\textbf{Correction pour comparaisons multiples.} Quatre tests $t$ appariés sont menés
sur ce même jeu de résultats ; une correction de Bonferroni ($\alpha = 0{,}05/4 =
0{,}0125$) est donc appliquée par prudence. \textbf{Les deux effets de durée et la
granularité à 8 époques restent significatifs après correction} (durée, réseau entier,
$p=1{,}3\times10^{-3}$ ; durée, tête seule, $p=1{,}2\times10^{-2}$ ; granularité,
8 époques, $p=6{,}2\times10^{-3}$, un ordre de grandeur sous le seuil ajusté) ; \textbf{
seule la granularité à 30 époques ne survit plus au seuil corrigé}
($p=2{,}2\times10^{-2}$). Ceci est cohérent avec le reste de la section plutôt qu'une
anomalie isolée : à 30 époques, le réseau entier a eu plus de temps pour rattraper une
partie de son retard grâce à l'effet de durée démontré ci-dessus, ce qui resserre
mécaniquement l'écart entre les deux granularités et fragilise statistiquement
précisément cette comparaison-là. Le message principal n'en est pas affaibli : c'est la
comparaison la moins centrale des deux (celle où l'écart se resserre déjà), pas celle
qui porte la conclusion.

\textbf{À granularité égale, personnaliser plus longtemps aide un peu ; à durée égale,
personnaliser seulement la tête aide beaucoup plus.} La tête seule limite probablement
le surapprentissage sur le peu de données locales disponibles par client (34
paramètres ajustés contre 1\,362 pour le réseau entier) -- une explication plausible,
cohérente avec les pratiques usuelles de fine-tuning sur domaine cible réduit, mais non
vérifiée plus avant ici.

Ce résultat révise la lecture de la Section~\ref{subsec:fedper} : l'avantage de FedPer
vient à la fois du fait qu'une personnalisation ait lieu et, dans une bien moindre
mesure, de sa durée continue -- mais surtout de sa \emph{granularité} (tête seule), le
facteur dominant, un choix de conception qui s'avère pertinent et pas un simple détail
d'implémentation.

\paragraph{Limite à ne pas occulter.} L'absence de modèle global unique a un coût côté
déploiement : il n'existe pas de firmware universel à distribuer à toutes les
infirmières -- chaque wearable doit recevoir (et mettre à jour) sa propre tête de
classification en plus de l'extracteur partagé, ce qui complique légèrement la logistique
de déploiement par rapport à FedProx. Le gain de communication mesuré ($-2{,}5\,\%$) est
réel mais mécaniquement modeste sur une architecture déjà minuscule ($1\,362$
paramètres) : la tête ne pèse que 34 paramètres sur ce total.

\subsection{Explorer la meilleure architecture possible hors contrainte TinyML}
\label{subsec:fedrep_scaffold}

\textbf{Cette sous-section relâche volontairement la contrainte TinyML} (empreinte
microcontrôleur, taille de firmware) pour explorer la meilleure architecture atteignable
sous contraintes de \emph{sécurité seules} (pas de données brutes centralisées, garantie
DP formelle sur ce qui est transmis) -- elle n'est pas comparable aux sections
précédentes en termes d'objectif de déploiement embarqué, et ne remplace pas
l'architecture MLP Tiny déployable qui reste la contribution principale de ce travail.

\paragraph{Architecture.} Trois éléments distinguent cette architecture de FedPer : (1)~une
\textbf{tête de classification à deux couches} (16$\to$16$\to$2, 306 paramètres) au lieu
d'une seule couche linéaire (34 paramètres), toujours jamais transmise ni agrégée ;
(2)~le \textbf{DP-SGD restreint à l'extracteur partagé} -- le clipping et le bruit
gaussien d'Opacus ne s'appliquent qu'aux paramètres réellement transmis au serveur, la
tête étant entraînée sans bruit ; (3)~\textbf{SCAFFOLD}~\cite{karimireddy2020scaffold}
remplace le terme proximal de FedProx comme mécanisme de correction de la dérive
Non-IID sur l'extracteur partagé : chaque client maintient une variable de contrôle
$c_i$, transmise au serveur sous une clé préfixée dédiée pour bénéficier de
l'agrégation native, avec une mise à l'échelle côté client garantissant que la moyenne
pondérée par taille de partition (standard pour les poids du modèle) reconstruit bien la
moyenne simple non pondérée que SCAFFOLD exige pour $c_{\text{global}}$.

\textbf{Deux facteurs changent simultanément par rapport à FedPer, non démêlés dans ce
travail} : la capacité de la tête (1 couche contre 2) et le mécanisme de correction de
dérive (FedProx contre SCAFFOLD). Un résultat meilleur que FedPer ne permet pas de dire
lequel des deux domine, sauf ablation supplémentaire non réalisée ici -- même limite
méthodologique que celle déjà posée entre FedPer et FedProx (Section~\ref{subsec:fedper}).

\paragraph{Un incident numérique corrigé en cours de route, documenté par transparence.}
Le premier passage de la grille a divergé (NaN) : la correction SCAFFOLD, appliquée
après le clipping DP-SGD, n'était elle-même jamais bornée, alors que le gradient réel
l'est (à $\text{max\_grad\_norm}$). $c_i$ s'accumulant round après round, la correction
a fini par dominer le gradient protégé -- une boucle de rétroaction qui a fait exploser
les poids en quelques rounds. Corrigé en bornant la norme de la correction à la même
échelle que le clipping DP avant de l'ajouter au gradient, validé sur un balayage étendu
(12 rounds, pas seulement les 3 rounds du test rapide initial) avant de rejouer la
grille complète -- 25/25 runs, aucun NaN, aucun échec.

\paragraph{Découverte méthodologique : la domination par les clients mono-classe.} En
creusant un résultat de balanced-accuracy en apparence remarquablement stable
(quasi identique à deux décimales près sur plusieurs graines), l'inspection par client a
révélé que \textbf{11 des 15 clients ont un test-set mono-classe} (une seule classe
présente, split déterministe donc identique pour toute graine, Section~\ref{subsec:dataset}) :
une prédiction constante y obtient mécaniquement $100\,\%$, indépendamment de la qualité
réelle du modèle. Ces 11 clients dominent la balanced-accuracy groupée sur les 15
clients, masquant la performance réelle sur les 4 clients au test-set effectivement mixte
(clients 4, 5, 8, 12). \textbf{Cette limite s'applique rétroactivement à toutes les
métriques groupées sur 15 clients rapportées plus haut dans ce document}
(Tableaux~\ref{tab:main_results}, \ref{tab:dp_decomposition}, \ref{tab:fedper_comparison}) --
signalé explicitement en synthèse (Section~\ref{subsec:synthesis_limits}) plutôt que
laissé implicite. Une seconde métrique, restreinte aux 4 clients informatifs, est donc
calculée en complément pour les trois architectures.

\begin{table}[h!]
  \centering
  \caption{Comparaison FedProx+DP / FedPer / FedRep+SCAFFOLD (moyenne sur
  5$\sigma\times$5 graines, n=25 chacun) -- pooled sur 15 clients vs restreint aux 4
  clients au test-set réellement mixte.}
  \label{tab:fedrep_scaffold_comparison}
  \begin{tabular}{lcccc}
    \toprule
    Balanced-accuracy & FedProx+DP & FedPer & FedRep+SCAFFOLD \\
    \midrule
    Pooled (15 clients)            & $47{,}86 \pm 0{,}35\,\%$ & $65{,}06 \pm 2{,}78\,\%$ & $\mathbf{77{,}57 \pm 0{,}72\,\%}$ \\
    Clients informatifs (4/15)     & $42{,}73 \pm 0{,}97\,\%$ & $48{,}44 \pm 3{,}65\,\%$ & $\mathbf{59{,}89 \pm 0{,}22\,\%}$ \\
    AUC-ROC (clients informatifs)  & $0{,}2047 \pm 0{,}0232$  & $0{,}4417 \pm 0{,}1022$  & $\mathbf{0{,}7269 \pm 0{,}0035}$ \\
    \bottomrule
  \end{tabular}
\end{table}

Sur les clients informatifs -- la mesure la plus honnête cliniquement -- la hiérarchie
FedProx+DP $<$ FedPer $<$ FedRep+SCAFFOLD est statistiquement nette (Welch et
Mann-Whitney sur les 3 comparaisons deux-à-deux, $p < 10^{-7}$ dans tous les cas, détail
complet dans \texttt{resultsfeat/fedper\_results\_summary.md}). Fait notable, l'AUC de
FedProx+DP tombe à $0{,}20$ sur ces clients (pire que le hasard), confirmant que le
problème d'agrégation identifié Section~\ref{subsec:why_cent_wins} est précisément
concentré sur les clients les plus hétérogènes. FedRep+SCAFFOLD est aussi
\textbf{nettement plus stable} que les deux autres conditions sur ces clients
(écart-type de balanced-accuracy $0{,}22\,\%$ contre $3{,}65\,\%$ pour FedPer et
$0{,}97\,\%$ pour FedProx+DP) -- vérifié non-artefactuel (les checkpoints de graines
différentes ont été comparés directement et diffèrent réellement, pas de recyclage de
fichier).

\paragraph{Limite : coût de calcul non négligeable.} Cette grille a nécessité $\approx$
8h30 de calcul (deux passages complets suite à l'incident ci-dessus), contre
$\approx$ 1h40 pour FedProx+DP/FedPer -- un facteur $\approx 5\times$ plus lent par run,
plausiblement dû à la boucle de correction SCAFFOLD non vectorisée (itération Python sur
les paramètres à chaque pas local, jusqu'à $\approx 1240$ pas par round pour les plus
gros clients). Non profilé rigoureusement, mais un coût de calcul réel à ne pas ignorer,
même en dehors de toute contrainte de taille de modèle.

\subsection{Pourquoi le compromis confidentialité vs utilité est-il si plat ?}

Nos résultats montrent une balanced-accuracy quasiment constante sur toute la plage de
bruit testée (de $47{,}87\,\%$ à $\sigma=0{,}6$ à $48{,}12\,\%$ à $\sigma=1{,}8$, un
écart plus petit que la variance inter-graines), alors que $\varepsilon$ varie d'un
facteur $\approx 10$ sur cette même plage (Tableau~\ref{tab:epsilon_grid_extended}).
Deux éléments expliquent cette platitude :

\begin{itemize}
  \item L'Ablation de personnalisation (Section~\ref{subsec:ablation_personalization})
  montre que le modèle global lui-même reste stable en fonction de $\sigma$
  ($77$--$78\,\%$ de précision interne quel que soit le niveau de bruit) : l'essentiel
  de la dégradation liée à la confidentialité est donc déjà consommé à
  $\sigma = 0{,}6$, et l'augmenter encore n'ajoute pas de coût supplémentaire mesurable
  dans la plage testée.
  \item Le clipping de gradient ($C=1{,}0$) borne la contribution de chaque exemple
  indépendamment du niveau de bruit ajouté ensuite -- c'est cette borne, plus que
  l'amplitude du bruit lui-même, qui semble dominer la dynamique d'apprentissage sur
  ce modèle de petite taille et ce jeu de features.
\end{itemize}

Ce constat n'a pas pu être vérifié à l'échelle d'un modèle plus grand sur les données
corrigées (l'ablation de taille de modèle nécessite d'être rejouée, voir
Section~\ref{subsec:synthesis}) et ne doit donc pas être généralisé au-delà de
l'architecture MLP Tiny testée ici.

\subsection{Pourquoi la quantification INT8 s'effectue-t-elle presque sans perte de
précision en FL ?}

En mode centralisé, la quantification dynamique post-entraînement dégrade la précision
interne de $0{,}52 \pm 0{,}65\,\%$, tandis qu'en mode FL + DP, la dégradation mesurée
est quasi nulle et même légèrement négative en moyenne ($\approx -0{,}05\,\%$, à
l'intérieur du bruit de mesure), c'est-à-dire que la version quantifiée INT8 est
parfois marginalement plus précise que la version FP32 dont elle est dérivée.

Ce phénomène s'explique par la mécanique interne de DP-SGD~\cite{abadi2016deep}. Pour
borner la sensibilité des données, DP-SGD applique un \textbf{clipping de gradient}
strict ($C=1{,}0$) avant d'injecter le bruit, ce qui restreint la vitesse d'évolution
des paramètres et agit comme une régularisation implicite de type Tikhonov
(\textit{weight decay}).

Comme le montre le Tableau~\ref{tab:dynamic_range}, cette régularisation compresse
fortement la plage dynamique des poids par rapport au centralisé. À $\sigma = 0{,}8$,
la plage dynamique globale du modèle FL + DP est de $3{,}26 \pm 0{,}47$ (contre
$9{,}45 \pm 1{,}03$ pour le centralisé, soit un facteur $\approx 3$), et tombe à
$0{,}96 \pm 0{,}04$ sur la première couche (contre $7{,}11 \pm 0{,}57$). Les poids du
modèle FL étant nettement plus resserrés autour de zéro, ils se répartissent de
manière plus uniforme sur la grille de quantification INT8 à 256 niveaux, subissant
moins d'effet de saturation ou de perte de résolution que ceux du modèle centralisé,
dont la plage dynamique plus large est davantage tronquée par la quantification.

\subsection{Pourquoi la consommation de RAM pic varie-t-elle avec le bruit ?}

La fluctuation de la RAM pic observée lors de l'appel à \texttt{get\_epsilon()}
(de $11{,}34$~Mo à $\sigma=1{,}8$ jusqu'à $33{,}65$~Mo à $\sigma=0{,}6$) est une
conséquence directe de l'algorithme comptable d'Opacus~\cite{opacus}. En l'absence
de cet appel, l'empreinte mémoire propre au round d'entraînement local reste stable et
faible ($\approx 0{,}18$~Mo), indépendamment du niveau de bruit -- les objets Opacus
(moteur, optimiseur, \textit{dataloader}) étant instanciés une seule fois par client
au premier round puis réutilisés (Section~\ref{subsec:dp}), l'essentiel de
l'allocation mémoire liée au chargement des données a déjà eu lieu avant le round 30
et n'apparaît donc plus dans ce pic mesuré en fin de grille.

Pour calculer le budget $\varepsilon$, Opacus utilise par défaut l'approche PRV
(\textit{Privacy Random Variable}). Cet algorithme effectue une évaluation numérique
de la distribution de perte de confidentialité (PLD) accumulée sur l'historique des
gradients. Plus le multiplicateur de bruit $\sigma$ est faible, plus la distribution
est concentrée. Pour maintenir la précision numérique lors de l'intégration,
l'algorithme PRV génère une grille discrète beaucoup plus dense, proportionnelle à
$1/\sigma$. C'est cette grille numérique temporaire qui provoque le pic de RAM observé.

\subsection{Synthèse et perspectives}
\label{subsec:synthesis}

\subsubsection{Bilan}

Ce travail, appuyé sur une recherche sur grille systématique couvrant 25 runs fédérés
($\sigma \in [0{,}6;\, 1{,}8]$ sur 5 graines chacun) et plusieurs ablations ciblées,
aboutit à un résultat plus fondamental que la comparaison FL-vs-centralisé qui en était
le point de départ : \textbf{(0)~l'hétérogénéité extrême entre infirmières limite la
généralisation de toute architecture testée, centralisée comprise}. 11 des 15 clients
ont un test-set mono-classe (Section~\ref{subsec:dataset}), qui gonfle mécaniquement
toute métrique groupée sur 15 clients ; restreinte aux 4 clients au test-set réellement
mixte, la mesure la plus honnête cliniquement, le centralisé lui-même s'effondre
(balanced-accuracy $45{,}73\,\%$, AUC $0{,}267$, sous le hasard,
Section~\ref{subsec:why_cent_wins}) -- une limite de fond partagée par toutes les
architectures, pas un problème propre au FL. \textbf{Portée à garder à l'esprit} : la
significativité statistique de cet écart ($p<10^{-4}$ par endroits) est reproductible
d'une graine et d'un niveau de bruit à l'autre, mais porte structurellement sur
seulement 4 clients au test-set mixte -- ce jeu de données n'en offre pas d'autre pour
vérifier la généralisation à un 5\textsuperscript{e} profil clinique
(Section~\ref{subsec:why_cent_wins}).

Sur cette base, ce travail montre qu'il est possible de construire un système FL
opérationnel qui : (1)~fournit une garantie formelle $(\varepsilon,
\delta)$-DP~\cite{dwork2014foundations} atteignant $\varepsilon = 0{,}54 \pm 0{,}04$ à
$\sigma = 1{,}8$ (sous $\varepsilon = 1$, un seuil couramment retenu dans la littérature
DP pour des applications sensibles) ; (2)~quantifie honnêtement, sur les deux
métriques (pooled et clients informatifs), le coût de cette confidentialité et de cette
décentralisation par rapport à une baseline centralisée non déployable en pratique sur
ce type de données : un écart de 10 points de balanced-accuracy sur la métrique pooled
(57,84~\% contre 47,9~\%) qui se réduit à 3 points sur les clients informatifs
(45,73~\% contre 42,73~\%), pour une précision brute statistiquement équivalente dans
les deux cas (82,0~\% contre 81,95~\%) ; (2bis)~décompose ce coût grâce à une ablation
FedProx sans DP-SGD (Tableau~\ref{tab:dp_decomposition}) et montre qu'il est attribuable
essentiellement à la décentralisation elle-même (FL sans DP : $48{,}68\,\%$ pooled /
$45{,}18\,\%$ clients informatifs, contre $47{,}83\,\%$ / $42{,}73\,\%$ pour FL+DP), le
bruit de confidentialité n'ajoutant qu'un coût secondaire, plus visible sur les clients
difficiles (2,45 points) que sur la métrique pooled (moins d'un point) ; (3)~produit un
modèle déployable sur microcontrôleur en 5,32~Ko (FP32), quantifiable dynamiquement en
INT8 avec une dégradation de précision quasi nulle, expliquée par le clipping de
gradient DP-SGD qui resserre fortement la plage dynamique des poids ; (4)~présente un
compromis confidentialité/utilité empiriquement plat sur toute la plage testée
$\varepsilon \in [0{,}54;\, 5{,}29]$ ; (5)~teste une architecture qui supprime
l'agrégation de la couche de classification (FedPer, Section~\ref{subsec:fedper}) : sur
les clients informatifs, la balanced-accuracy remonte à $48{,}44 \pm 3{,}65\,\%$ (contre
$42{,}73\,\%$ pour FedProx+DP standard, +5,7 points, statistiquement significatif mais
d'ampleur modeste comparé au gonflement apparent sur la métrique pooled, +17,5 points) --
au prix d'une contrainte de déploiement propre à cette architecture (pas de modèle
global unique, chaque client garde sa propre tête de classification) ; et (6)~pousse
cette logique plus loin hors contrainte TinyML (Section~\ref{subsec:fedrep_scaffold}) en
combinant tête multi-couches, DP restreint à l'extracteur et SCAFFOLD : sur les clients
informatifs, la hiérarchie FedProx+DP ($42{,}7\,\%$) $<$ FedPer ($48{,}4\,\%$) $<$
FedRep+SCAFFOLD ($59{,}9\,\%$, aussi nettement plus stable) est statistiquement nette --
au prix d'un coût de calcul $\approx 5\times$ supérieur, non négligeable même hors
contrainte de taille de modèle. Dans tous les cas, ces architectures de
personnalisation réduisent l'écart sans jamais le faire disparaître : même
FedRep+SCAFFOLD reste loin d'un pouvoir discriminant pleinement satisfaisant sur les
clients les plus atypiques, cohérent avec le constat (0) que le problème de fond est
l'hétérogénéité elle-même, pas seulement le choix d'architecture FL.

\textbf{Le résultat le plus directement actionnable de ce travail, cependant, ne vient
d'aucune des architectures ci-dessus.} L'isolation du facteur granularité
(Section~\ref{subsec:duration_vs_granularity}) montre qu'un simple fine-tuning hors
ligne de la seule tête de classification, pendant 8 époques, sur le checkpoint
FedProx+DP \emph{déjà existant} et sans aucune modification du pipeline
d'entraînement fédéré, atteint $55{,}88 \pm 2{,}35\,\%$ de balanced-accuracy sur les
clients informatifs -- \textbf{davantage que FedPer} ($48{,}44 \pm 3{,}65\,\%$), une
architecture qui nécessite pourtant un chemin d'entraînement fédéré dédié, une
sauvegarde de tête par client et une méthodologie d'évaluation propre. Pour un
déploiement réel, ceci suggère qu'il n'est probablement pas nécessaire de complexifier
l'architecture fédérée elle-même pour obtenir l'essentiel du gain de personnalisation :
un post-traitement léger et peu coûteux à intégrer (geler l'extracteur, fine-tuner la
tête sur les données locales après réception du modèle global standard) suffit, avec
une charge d'ingénierie et de déploiement nettement moindre que celle de FedPer ou de
FedRep+SCAFFOLD.

\subsubsection{Limites et travaux futurs}
\label{subsec:synthesis_limits}

Plusieurs limites méritent d'être reportées honnêtement :

\begin{itemize}
  \item \textbf{Ablation de taille de modèle à rejouer.} L'ablation comparant MLP Tiny
  et MLP Medium n'a pas encore été rejouée sur le pipeline actuel et n'est donc pas
  rapportée dans ce travail.
  \item \textbf{$\delta$ légèrement supérieur au seuil recommandé.} Avec
  l'élargissement du jeu d'entraînement à $n = 306\,886$ fenêtres après recombinaison
  des labels, $\delta = 10^{-5}$ dépasse désormais d'un facteur $\approx 3$ le seuil
  usuel $\delta < 1/n$ (Tableau~\ref{tab:dp_params_justif}). Les valeurs
  d'$\varepsilon$ rapportées restent correctement calculées pour ce $\delta$, mais un
  $\delta = 10^{-6}$ serait recommandé pour une itération future.
  \item \textbf{Mesure de la taille du modèle quantifié non fiable.} La fonction de
  calcul de taille de modèle (\texttt{get\_model\_size}) énumère les paramètres via
  \texttt{model.parameters()}, qui ne couvre pas les poids INT8 empaquetés par
  \texttt{torch.quantization.quantize\_dynamic}. La taille INT8 effective n'a donc pas
  pu être mesurée de façon fiable dans ce travail ; seul le facteur de compression
  théorique ($\times 3{,}8$, Section~\ref{subsec:quantization}) est rapporté.
  \item \textbf{Mesure RAM liée à l'implémentation Opacus.} L'empreinte mémoire
  pendant le calcul d'$\varepsilon$ dépend fortement de l'implémentation interne du
  PRV accountant~\cite{opacus}. Une future étude pourrait implémenter un comptable
  analytique (basé sur le \textit{moments accountant} d'origine~\cite{abadi2016deep})
  pour éviter cette surcharge numérique en environnement très contraint.
  \item \textbf{Simulation sur nœud unique.} Notre évaluation a été exécutée sous
  forme de simulation via Flower~\cite{beutel2020flower} sur un seul nœud de calcul.
  Valider les performances réelles nécessiterait un déploiement décentralisé réel sur
  des wearables physiques communicant par canaux sécurisés.
  \item \textbf{Coût de la décentralisation non encore réduit.} La seule mitigation
  testée en agrégation globale (loss pondérée par classe, Section~\ref{subsec:why_cent_wins})
  n'a pas réduit l'écart de balanced-accuracy avec le centralisé. D'autres pistes --
  features physiologiques plus informatives, repondération de la contribution de
  chaque client dans l'agrégation FedProx, personnalisation renforcée -- n'ont pas été
  explorées dans ce travail et constituent la suite naturelle de cette étude. Une piste
  plus ciblée, et la mieux étayée empiriquement parmi celles listées ici, ressort du
  scénario du nouvel arrivant (Section~\ref{subsec:late_joiner_recall}) : le seuil de
  décision optimal après adaptation locale diverge en sens opposé selon le client
  ($0{,}001$ contre $0{,}950$ sur les deux cas testés), ce qui suggère qu'une
  calibration de seuil \emph{par client}, réalisée sur un split de validation dédié et
  distinct du test, pourrait apporter un gain de rappel clinique sans modifier le
  modèle lui-même -- une piste peu coûteuse à valider dans une itération future.
  \item \textbf{Résultats spécifiques à ce jeu de données et cette configuration.}
  Le coût mesuré de la confidentialité, l'erreur de quantification quasi nulle et la
  platitude du compromis confidentialité/utilité sont des observations empiriques
  propres à ce jeu de données Non-IID de stress infirmier et à cette petite
  architecture MLP. Elles ne doivent pas être généralisées sans validation sur
  d'autres jeux de données et tailles de modèles.
  \item \textbf{FedPer vs FedProx, durée et granularité : isolés séparément, à partir
  du même checkpoint, avec test statistique apparié.} L'écart initialement observé
  entre FedPer et FedProx fine-tuné (Section~\ref{subsec:fedper}) aurait pu refléter
  soit la \emph{durée} de personnalisation, soit la \emph{granularité}. Isolés l'un
  après l'autre (Section~\ref{subsec:duration_vs_granularity}), en partant à chaque
  fois du même checkpoint FedProx+DP et avec un fine-tuning re-seedé (5 graines,
  test $t$ apparié) : la granularité domine largement (+6,5 à +7,4 points,
  $p<0{,}03$ -- solidement établi à 8 époques, $p=6{,}2\times10^{-3}$, encore un ordre
  de grandeur sous le seuil de Bonferroni ajusté pour les 4 tests menés,
  $\alpha=0{,}0125$) -- personnaliser uniquement la tête surpasse nettement le
  fine-tuning du réseau entier, à extracteur de départ identique -- et la durée compte
  aussi, plus modestement (+1 à +2 points, $p<0{,}02$, les deux valeurs survivant à
  cette même correction). Deux corrections successives ont été
  nécessaires pour arriver à ce résultat : un premier test comparant directement
  FedPer (son propre extracteur) au fine-tuning complet de FedProx confondait la
  granularité avec le fait que FedPer part d'un extracteur entraîné selon un régime
  différent ; un deuxième test, corrigeant ce point mais sans re-seeder le
  fine-tuning, avait conclu à tort que la durée n'avait aucun effet -- le bruit
  d'ordre de mélange non contrôlé masquait un effet réel mais petit. Le test de
  Wilcoxon apparié plafonne à $p=0{,}0625$ sur les quatre comparaisons, un artefact de
  plancher propre à $n=5$ paires de même signe, pas une significativité douteuse.
  \item \textbf{Domination des métriques groupées par les clients mono-classe --
  corrigé dans ce document, pas seulement documenté.} Découvert lors de l'évaluation de
  FedRep+SCAFFOLD (Section~\ref{subsec:fedrep_scaffold}) : 11 des 15 clients ont un
  test-set mono-classe (Section~\ref{subsec:dataset}), où une prédiction constante
  obtient mécaniquement $100\,\%$ quelle que soit la qualité du modèle. Une
  vérification systématique a été menée sur \emph{tous} les tableaux pooled de ce
  document (Tableaux~\ref{tab:main_results}, \ref{tab:dp_decomposition},
  \ref{tab:fedper_comparison}, \ref{tab:fedrep_scaffold_comparison}) : chacun rapporte
  désormais la métrique pooled sur 15 clients \emph{et} la métrique restreinte aux 4
  clients au test-set réellement mixte (4, 5, 8, 12), la seconde étant la mesure la
  plus honnête cliniquement. Les écarts relatifs entre conditions restent
  qualitativement dans le même sens sur les deux métriques, mais d'ampleur nettement
  réduite sur la métrique informative pour la comparaison centralisé-vs-FL (10 points
  $\to$ 3 points) et FedPer-vs-FedProx (17,5 points $\to$ 5,7 points) -- un changement
  de magnitude, pas seulement un ajustement cosmétique, qui a motivé la reformulation du
  message central de ce document (Section~\ref{subsec:why_cent_wins}).
\end{itemize}

\printbibliography

\end{document}